\documentclass[11pt]{article}

\DeclareUnicodeCharacter{0301}{\'e}

\usepackage[margin=1in]{geometry}
\usepackage{multirow}
\usepackage{amssymb,amsmath,amsthm,amsfonts}
\usepackage{bm}
\usepackage{textgreek}
\usepackage{mathtools}
\usepackage{enumitem}
\usepackage[numbers,comma,sort&compress]{natbib}
\usepackage{authblk}
\usepackage{graphicx}
\usepackage[font=small]{caption}
\usepackage{subcaption} % [labelformat=simple]
\usepackage{tablefootnote} 
\usepackage{float}
\usepackage[ruled,vlined,linesnumbered]{algorithm2e}
\usepackage{algorithmic}
\renewcommand{\algorithmicrequire}{\textbf{Input: }}  
\renewcommand{\algorithmicensure}{\textbf{Output: }}
\usepackage{physics}


\usepackage{footnote}
\usepackage{xcolor}
\usepackage{mathrsfs}
\usepackage{bbm}
\usepackage{makecell}
\usepackage{pbox}
\usepackage[colorlinks]{hyperref}
\hypersetup{citecolor=blue}
% \urlstyle{same}
%\captionsetup[figure]{labelfont=bf}

\newtheorem{fact}{Fact}[section] 
\newtheorem{theorem}{Theorem}
\newtheorem{definition}{Definition}
\newtheorem{lemma}{Lemma}
\newtheorem{proposition}{Proposition}
\newtheorem{example}{Example}
\newtheorem{corollary}{Corollary}
\newtheorem{remark}{Remark}
\newtheorem{assumption}{Assumption}
\newcommand{\eqn}[1]{(\ref{eqn:#1})}
\newcommand{\eq}[1]{(\ref{eq:#1})}
\newcommand{\prb}[1]{(\ref{prb:#1})}
\newcommand{\thm}[1]{\hyperref[thm:#1]{Theorem~\ref*{thm:#1}}}
\newcommand{\cor}[1]{\hyperref[cor:#1]{Corollary~\ref*{cor:#1}}}
\newcommand{\defn}[1]{\hyperref[defn:#1]{Definition~\ref*{defn:#1}}}
\newcommand{\lem}[1]{\hyperref[lem:#1]{Lemma~\ref*{lem:#1}}}
\newcommand{\prop}[1]{\hyperref[prop:#1]{Proposition~\ref*{prop:#1}}}
\newcommand{\assum}[1]{\hyperref[assum:#1]{Assumption~\ref*{assum:#1}}}
\newcommand{\fig}[1]{\hyperref[fig:#1]{Figure~\ref*{fig:#1}}}
\newcommand{\tab}[1]{\hyperref[tab:#1]{Table~\ref*{tab:#1}}}
\newcommand{\algo}[1]{\hyperref[algo:#1]{Algorithm~\ref*{algo:#1}}}
\renewcommand{\sec}[1]{\hyperref[sec:#1]{Section~\ref*{sec:#1}}}
\newcommand{\append}[1]{\hyperref[append:#1]{Appendix~\ref*{append:#1}}}
\newcommand{\fac}[1]{\hyperref[fac:#1]{Fact~\ref*{fac:#1}}}
\newcommand{\lin}[1]{\hyperref[lin:#1]{Line~\ref*{lin:#1}}}
\newcommand{\itm}[1]{\ref{itm:#1}}
\newcommand{\bdelta}{\mathbf{\Delta}}
\renewcommand{\arraystretch}{1.25}
\newcommand{\algname}[1]{\textup{\texttt{#1}}}
\renewcommand{\i}{\mathrm{i}}
\renewcommand{\d}{\mathrm{d}}
\renewcommand{\L}{\mathcal{L}}
\newcommand{\ad}{\mathrm{ad}}
\newcommand{\T}{\mathcal{T}}
\newcommand{\B}{\Big}
\renewcommand{\P}{\mathcal{P}}
\newcommand{\simu}{\mathrm{sim}}
\newcommand{\Hc}{H_{\mathrm{control}}}
\newcommand{\He}{H_{\mathrm{eff}}}
\newcommand{\A}{\mathcal{A}}
\renewcommand{\O}{\mathcal{O}}
\newcommand{\tot}{\mathrm{tot}}
\newcommand{\G}{\mathcal{G}}
\newcommand{\E}{\mathcal{E}}
\renewcommand{\comm}{\mathrm{comm}}
\newcommand{\supp}{\mathrm{supp}}
\newcommand{\eps}{\varepsilon}
\newcommand{\mc}{\mathcal}
\newcommand{\mb}{\mathbb}
\def\>{\rangle}
\def\<{\langle}
\def\trans{^{\top}}

\newcommand{\xz}[1]{{\color{cyan} XZ: #1}}
\newcommand{\tyl}[1]{{\color{blue} Tongyang: #1}}
\newcommand{\syz}[1]{{\color{teal} (Shengyu: #1)}}
\newcommand{\syzn}[1]{{\color{teal} #1}}
\newcommand{\sz}[1]{{\color{blue} Shuo: #1}}

\allowdisplaybreaks

\title{Jointly Mitigating Physical and Trotter Errors in Hamiltonian Simulation with Commutator Bounds
}
\author{}
\date{}

\begin{document}

\maketitle

Given an $n$-qubit Hamiltonian $H = \sum_{\gamma=1}^{\Gamma} H_{\gamma}$, we aim to simulate $e^{-\i H T}$ by Trotterization under Markovian physical noise. 

\textcolor{red}{Yicong: In addition to the number of interpolation nodes, try to prove that the overall sample complexity of ZNE procedure via 1D extrapolation is smaller than that by 2D extrapolation.}

\section{Introduction}

\xz{I tend to think that jointly mitigate algorithmic and physical error is a minor contribution of our work. The main technique contribution of our work is an upper bound of the truncation error of the BCH formula for general Lindbladians in terms of the doubly right-nested commutators. Compared to \cite{mizuta2025commutator}, we go beyond lattice Hamiltonian to general cases. Compared to \cite{fang2025high}, we obtain fine-grained commutator bounds, and use a different way to bound the remainder to incorporate dissipative dynamic.} 
If we use this perspective, I think our abstract should contain like \emph{Exponentially Reduced Circuit Depth for Trotterization under Physical Noise with Rigorous (General) Commutator Bounds}
The storyline of intro can be:
\begin{itemize}
    \item Hamiltonian simulation is a promising problem to show quantum advantages on near-term devices, and list some applications. \xz{We need to emphasize near-term applications, which makes it natural to use Trotter and consider physical noise. }
    \item List some Hamiltonian simulation algorithms, and state that although Trotter method has sub-optimal dependence in $t$ and $\eps$, it has some advantages like commutator bounds, no ancilla and complex controlled gates, which makes it more practical on near-term devices. 
    \item \citet{watson2025exponentially} exponentially reduced the error dependence of Trotter using extrapolation. However, providing a fully rigorous commutator bounds of general Hamiltonians are challenging. A key technical difficulty is that the convergence of the BCH formula, which the analysis relies on, is not guaranteed in all regimes (e.g., for general $k$-local Hamiltonians), making the derivation of commutator bounds difficult. Important progress was recently made by \citet{mizuta2025commutator}, who provided a rigorous commutator bound for the specific case of lattice Hamiltonians by using a subsystem Trotterization technique. This naturally leaves open the crucial question: \begin{align*}
        \textit{Can we establish a commutator bound for Trotter with extrapolation for general Hamiltonians? }
    \end{align*}  Then we can state our first contribution: \emph{We develop a new analysis, centered on doubly right-nested commutators, that provides a rigorous error bound for general Hamiltonians.} We apply our result to $k$-local Hamiltonian on lattice and second quantized Hamiltonian. The main technique ingredient is an upper bound of the truncation error of the BCH formula for general Hamiltonians in terms of the doubly right-nested commutators.
    \item Near-term devices suffer from physical nose. The physical error can be reduced by qec or qem. qec is scalable but expensive, qem could incur an exponential sampling overhead in the circuit depth, but is more friendly to near-term devices. Introduce QEM methods, especially ZNE. 
    \item Most existing methods handle these two parts of error separately. Recent works show that the Trotter error and physical error can be simultaneously mitigated by a one-dimensional extrapolation. However, they only consider first order Trotter, model the physical noise by a global depolarizing channel approximately, and commutator bounds are missing. Then we ask the second question \begin{align*}
        \textit{Can we jointly mitigate physical and Trotter error while maintaining the commutator bound?}
    \end{align*}Our analysis of the truncation error of the BCH formula naturally applies to Lindbladian, which makes us to bound the extrapolation error of jointly mitigating dissipative physical error and algorithmic error in terms of the nested commutators of the Hamiltonians and noise Lindbladians. Then we can state our second contribution: \emph{We propose an algorithm that jointly mitigates physical and Trotter error with commutator bounds and provide the optimal choice of extrapolation parameters for some important Hamiltonians.} \xz{maybe add in discussion: Our analysis takes the locality of noise into account which can be used to give specific analysis of quantum devices with certain connectivity. }
\end{itemize}



\section{Notations}
For a sequence of operators $\A_1, \ldots, \A_{\Gamma}$, we denote $\prod_{\gamma=1}^\Gamma \A_{\gamma}= A_{\Gamma}\A_{\Gamma-1} \ldots \A_1$ and $\prod_{\gamma=\Gamma}^1 \A_{\gamma}= A_{1}\A_{2} \ldots \A_\Gamma$. We denote the right-nested commutator $[X_1, [X_2, [\cdots[X_{n-1}, X_n]\cdots]]]$ by $[X_1, X_2, \ldots, X_n]$. We denote $\{1, 2, \ldots, n\}$ by $[n]$ and the collection of all size-$m$ subsets of $[n]$ by $\binom{[n]}{m}$. We use $\supp(X)$ to denote the set of qubits that $X$ nontrivially acts on.  For a superoperator $\mc L$, we denote the diamond norm of $\mc L$ by $\|\mc L\|$.
Define the time-ordering operator $\mc T\{\cdot\}$ as \begin{align*}
    \mc T\big\{[O_1(t_1),O_2(t_2),\ldots, O_q(t_q) ]\big\} = [O_{\pi(1)}(t_{\pi(1)}),O_{\pi(2)}(t_{\pi(2)}),\ldots, O_{\pi(q)}(t_{\pi(q)})],
\end{align*}
where $\pi\colon [q]\to[q]$ is a permutation of $[q]$ such that \[
    t_{\pi(1)}\ge t_{\pi(2)}\ge \cdots \ge t_{\pi(q)}.
    \]
For any permutation $\sigma\colon [q]\to [q]$, define the reordering operator $\mc R_{\sigma}\{\cdot\}$ as \begin{align*}
    \mc R_{\sigma}\big\{[O_1(t_1),O_2(t_2),\ldots, O_q(t_q)]\big\} = [O_{\sigma(1)}(t_{\sigma(1)}),O_{\sigma(2)}(t_{\sigma(2)}),\ldots, O_{\sigma(q)}(t_{\sigma(q)})].
\end{align*}

\section{Preliminaries}
\subsection{Quantum error mitigation}

\subsection{Magnus expansion}
Consider a time-dependent linear differential equation \begin{align*}
    \frac{\d Y(t)}{\d t} = A(t) Y(t), \quad Y(0) = I.
\end{align*} The Magus expansion expresses $Y(t)$ as a time-independent \syz{dependent?} exponential $Y(t) = \exp(\Omega(t))$. The exponent $\Omega(t)$ admits a series expansion \begin{align*}
    \Omega(t) = \sum_{q = 1}^{\infty} \Omega_q(t).
\end{align*} 
The $p$-th order Magnus expansion is denoted by \begin{align}
    \Omega_{(p)} = \sum_{q=1}^{p} \Omega_q.
\end{align}
Each $\Omega_q(t)$ can be written as an integral of nested commutators of $A(t)$:
\begin{align}
\label{eq:magnus-formula}
    \Omega_q(t) = \sum_{\sigma\in S_q} c_{\sigma, q}  \int_{0}^t \d \tau_1 \int_{0}^{\tau_1}\d \tau_2 \cdots \int_{0}^{\tau_{q-1}}\d \tau_q \, \mc R_{\sigma}\big\{[A(\tau_1), A(\tau_2), \ldots, A(\tau_q)] \big\},
\end{align}
where $S_q$ denotes the set of permutations of $[q]$, \begin{align*}
    c_{\sigma, q}=\frac{1}{q^2}\frac{(-1)^{d_{\sigma}}}{\binom{q-1}{d_\sigma}} 
\end{align*} satisfying $|c_{\sigma, q}|\le 1/q^2$, and $d_{\sigma}$ denotes the number of $i\in[q-1]$ such that $\sigma(i)>\sigma(i+1)$. The derivative of $\Omega_q(t)$ is
\begin{align}
    \dot{\Omega}_q(t) = \sum_{\sigma\in S_q} c_{\sigma, q}  \int_{0}^{t}\d \tau_2 \int_0^{\tau_2} \d\tau_3 \cdots \int_{0}^{\tau_{q-1}}\d \tau_q \, \mc R_{\sigma}\big\{[A(t), A(\tau_2), \ldots, A(\tau_q)] \big\}.
\end{align}
The exponential $Y_{(q)}(t) :=\exp(\Omega_{(q)}(t))$ is the solution of the following differential equation: \begin{align*}
    \frac{\d Y_{(q)}(t)}{\d t} = A_{(q)}(t) Y_{(q)}(t), \quad Y(0) = I,
\end{align*}
where \begin{align*}
    A_{(q)}(t)  = \sum_{k = 0}^{\infty}\frac{1}{(k+1)!}\ad_{\Omega_{(q)}}^k(\dot\Omega_{(q)}).
\end{align*}
For any expression $L$ involving integrals of nested commutators of $A(t)$, the grade of $L$, denoted by $\rm gd(L)$, is the number of $A(t)$ in the commutator. \citet{fang2025high} showed that $A_{(q)}(t)-A(t)$ contains only terms with grade at least $q$. 
\begin{lemma}[{\cite[Proposition~3]{fang2025high}}]
\label{lem:grade-cond}
    For any $1\le k\le q$, the sum of all terms in $A_{(q)}(t) -A(t)$ with grade $k$ vanishes. 
\end{lemma} 
\begin{lemma}[{\cite[Lemma~4]{fang2025high}}]
\label{lem:grade-cond-q+1}
    The grade $q+1$ terms in $A_{(q)}(t) -A(t)$ match those in $-\dot\Omega_{q+1}(t)$.
\end{lemma}
\subsection{Richardson extrapolation}

\begin{lemma}[{see \cite[Lemma A.9]{wang2025randomized}}]
\label{lem:richardson-extrapolation}
    Suppose that function $f\colon \mb R \to \mathbb{R}$ can be decomposed as 
    \begin{align*}
        f(x) = P_m(x)+R_m(x),
    \end{align*} 
    where $P_m(x)$ is a degree-$(m-1)$ polynomial and $R_m(x)$ is the remainder. Given any $t > 0$ and $s\in (0,1)$, we can choose 
    \begin{align*}
        K = \max\Big\{\frac{m}{\pi}, \frac{2t}{s}\Big\}, \quad r_j =\left\lceil \frac{K}{\sin^2(\pi(2j-1)/8m)}\right\rceil,
    \end{align*}
    and
    \begin{align*}
        s_j = \frac{t}{r_j}, \quad b_j=\prod_{\ell \neq j} \frac{1}{1-r_\ell/r_j}
    \end{align*}
    so that $\|\mathbf{b}\|_1 \le C\log m$ for some absolute constant $C$, $r_j=O(\max\{m^3, m^2t/s\}/j^2)$, and $\max_{j}\{ s_j \}\le s$. Then define the $m$-term Richardson extrapolation \begin{align}
        \label{eq_richardson extrapolation}
        F^{(m)}(s)=\sum_{j=1}^m b_j f(s_j),
    \end{align}
    which satisfies \begin{align}
    \label{eq:error-F}
        |F^{(m)}(s)-f(0)|\le \|\mathbf{b}\|_1 \max_j|R_{m}(s_j)|.
    \end{align}
\end{lemma}

\section{Main results}
We consider a general $p$-th order product formula of the form
\begin{align*}
    P(t) = \prod_{v=1}^Me^{-\i t a_{v} H_{\gamma_v}},
\end{align*}
where the product is ordered from right to left. Here, $a_{v}$ are real coefficients, $M = \Upsilon\Gamma$ is the total number of unitary gates, and the sequence $\gamma_v$ specifies which Hamiltonian is applied at step $v$. Denote $a_{\max} = \max_{v\in [M]} |a_v|$. The corresponding ideal channel $\P(t)$ is a superoperator given by
\begin{align*}
    \P(t) = \prod_{v=1}^Me^{-\i t a_{v} \A_{\gamma_v}},
\end{align*}
where $\A_{\gamma}:=\ad_{H_{\gamma}}$ denotes the adjoint action of $H_{\gamma}$. 

Suppose that the system is subject to physical noise with a controllable strength $\lambda>0$. Specifically, after each layer of unitary gates, a noise channel $e^{\lambda \mc L_{v}}$ is applied, where $\mc L_{v}$ is a Lindbladian. The resulting noisy channel is
\begin{align*}
    \tilde{\P}(t,\lambda)=\prod_{v=1}^M\left( e^{\lambda \mc L_{v}} e^{-\i t a_{v} \A_{\gamma_v}} \right).
\end{align*}
Consider a time-dependent linear differential equation \begin{align}
   \frac{\d \mc Y(\tau)}{\d \tau} = \mc L(\tau) \mc Y(\tau), \quad \mc Y(0)= \mc I, \quad \mc L(\tau) = \begin{cases} -\i t a_1 \A_{\gamma_1} & \text{if } \tau \in [0, 1] \\ \lambda \mc L_1 & \text{if } \tau \in (1, 2] \\ \vdots & \vdots \\ -\i t a_M \A_{\gamma_M} & \text{if } \tau \in (2M-2, 2M-1] \\ \lambda \mc L_M & \text{if } \tau \in (2M-1, 2M] \end{cases}, \label{eq:df-eq}
\end{align}
such that \[
\mc Y(2M) = \tilde{\mc P}(t, \lambda).
\]
For ease of presentation, we define \begin{align}
    \mathcal{K}_j := \begin{cases}\mathcal{A}_{\gamma_v} & \text{if } j = 2v-1\\ \mathcal{L}_v & \text{if } j = 2v \end{cases}, \quad \tilde{\mathcal{K}}_j := \begin{cases}-\i ta_{v}\mathcal{A}_{\gamma_v} & \text{if } j = 2v-1\\ \lambda\mathcal{L}_v & \text{if } j = 2v \end{cases} \label{eq:def-K}
\end{align}
where $v = 1, 2,\ldots, M$.

Let 
\[ \Omega_{t,\lambda}(\tau) = \sum_{q=1}^{\infty}\Omega_{q,t,\lambda}(\tau)\]
be the Magnus expansion of the solution $\mc Y(\tau)$ of Eq.~\eq{df-eq}, which means \begin{align*}
    \mc Y(\tau) = \exp(\Omega_{t,\lambda}(\tau)),
\end{align*} and \begin{align*}
    \tilde{\mc P}(t,\lambda) =\mc Y(2M)= \exp(\Omega_{t,\lambda}(2M))
\end{align*} especially. 

The Magnus expansion gives a series expansion for the effective generator of $\tilde{\P}(t,\lambda)$:
\begin{align}
\label{eq:bch}
    \tilde{\P}(t,\lambda) = \exp\bigg( \sum_{q=1}^{\infty} \Omega_{q,t,\lambda}(2M) \bigg),
\end{align}
for sufficiently small $t$ and $\lambda$. %The first-order term $\Omega_{1,t,\lambda}(2M)$ is the sum of all elementary generators:
%\begin{align*}
%    \Omega_{1,t,\lambda}(2M) = \sum_{v=1}^M \left( -\i t a_v \A_{\gamma_v} + \lambda \mc L_v \right) = -\i t \ad_{H} + \lambda\sum_{v=1}^{M}\mc L_{v}.
%\end{align*} 
In the following lemma, we show that the Magnus expansion matches the BCH formula of $ \tilde{\P}(t,\lambda)$ and provide an upper bound of each term in the Magnus expansion in terms of right-nested commutators of $\mc K_j$.

\begin{lemma}
\label{lem:magnus-bch}
    Define \begin{align}
    \label{eq:alpha-q-q}
    \alpha_{\comm}^{(q_1, q_2)}=\sum_{(j_1, \ldots, j_q)\in S_{q_1, q_2}}\big\|[\mc K_{j_1}, \mc K_{j_2},\ldots,\mc K_{j_q}]\big\|,
    \end{align}
    where \begin{align}
    \label{eq:S-q-q}
        S_{q_1, q_2}:=\Bigl\{(j_1, \ldots, j_q)\in [2M]^q\Bigm| |\{k\in [q]\mid j_k \text{ is  odd}\}|=q_1,~
        |\{k\in [q]\mid j_k \text{ is  even}\}|=q_2\Bigr\}.
    \end{align}
    The $q$-th order term $\Omega_{q,t,\lambda}(2M)$ in the Magnus expansion of $\mc Y(2M)$ in Eq.~\eq{df-eq} is identical to the $q$-th order term in the BCH formula of $\tilde{\P}(t,\lambda)=\prod_{v=1}^M\left( e^{\lambda \mc L_{v}} e^{-\i t a_{v} \A_{\gamma_v}} \right)$. Moreover, $\Omega_{q,t,\lambda}(2M)$ can be written as \begin{align}
    \Omega_{q,t,\lambda}(2M)= \sum_{\substack{q_1, q_2 \ge 0\\q_1+q_2 = q}}\Omega_{q_1, q_2} t^{q_1}\lambda^{q_2},
    \end{align}
    such that \begin{align}
    \|\Omega_{q_1, q_2}\|\le \frac{a_{\max}^{q_1}}{q^2}\alpha_{\comm}^{(q_1, q_2)}.
    \end{align}
\end{lemma}
\begin{proof}
    By Eq.~\eq{magnus-formula}, $\Omega_{q,t,\lambda}(2M)$ can be written as \begin{align}
        &\Omega_{q,t,\lambda}(2M) \nonumber\\
        =~&\sum_{\sigma\in S_q} c_{\sigma, q}  \int_{0}^{2M} \d \tau_1 \int_{0}^{\tau_1}\d \tau_2 \cdots \int_{0}^{\tau_{q-1}}\d \tau_q ~\mc R_{\sigma}\big\{[\mc L(\tau_1), \mc L(\tau_2), \ldots, \mc L(\tau_q)] \big\} \nonumber \\
        =~& \frac{1}{ q!}\sum_{\sigma\in S_q} c_{\sigma, q}  \int_{0}^{2M} \d \tau_1 \int_{0}^{2M}\d \tau_2 \cdots \int_{0}^{2M}\d \tau_q ~ \mc R_{\sigma} \circ \mc T\big\{[\mc L(\tau_1), \mc L(\tau_2), \ldots, \mc L(\tau_q)]\big\}  \nonumber\\
        =~&\frac{1}{ q!}\sum_{\sigma\in S_q} c_{\sigma, q} \sum_{ 1\le v_1,v_2,\ldots, v_q \le 2M} \int_{v_1-1}^{v_1} \d \tau_1 \int_{v_2-1}^{v_2}\d \tau_2 \cdots \int_{v_q-1}^{v_q}\d \tau_q ~ \mc R_{\sigma} \circ \mc T\big\{[\mc L(\tau_1), \mc L(\tau_2), \ldots, \mc L(\tau_q)] \big\}  \nonumber\\
        =~& \frac{1}{ q!}\sum_{\sigma\in S_q} c_{\sigma, q} \sum_{ 1\le v_1,v_2,\ldots, v_q \le 2M} \mc R_{\sigma} \circ \mc T\big\{[\mc L(v_1), \mc L(v_2), \ldots, \mc L(v_q)] \big\}, \label{eq:Z-q-first}
    \end{align}
    where the fourth line follows from $\mc L(\tau)$ is constant in $(v_j-1, v_j]$ for all $v_j = 1, 2, \ldots, 2M$. For any sequence of $v_1, v_2, \ldots, v_q$ in the summation in Eq.~\eq{Z-q-first} and $v\in [M]$, we denote the number of $v_j$ such that $\mc L(v_j) = \lambda \mc L_v$ by $p_v$, and the number of $v_j$ such that $\mc L(v_j) = -\i t a_v\mc A_{\gamma_v}$ by $p'_v$. Regrouping the summation in  Eq.~\eq{Z-q-first} gives \begin{align}
        &\Omega_{q,t,\lambda}(2M) \nonumber\\
        =~&\begin{multlined}[t] \frac{1}{ q!}\sum_{\sigma\in S_q} c_{\sigma, q} \sum_{\substack{p_v,p'_v\ge 0\nonumber\\ p_1+p_1'+\cdots+p_M+p_M'=q }}\binom{q}{p_1, p_1', \ldots, p_M, p_M'} \\ \cdot\mc R_{\sigma}\big\{[\underbrace{\lambda\mc L_M, \ldots, \lambda\mc L_M}_{p_M}, \underbrace{-\i ta_M\A_{\gamma_M}, \cdots, -\i ta_M\mc A_{\gamma_M}}_{p_M'},  \ldots, \underbrace{\lambda \mc L_1,  \ldots, \lambda\mc L_1}_{p_1}, \underbrace{-\i ta_1\A_{\gamma_1},  \ldots, -\i ta_1 \mc A_{\gamma_1}}_{p_1'}]\big\}\nonumber\end{multlined}\\
        =~& \begin{multlined}[t]
            \sum_{\substack{p_v,p'_v\ge 0\\ p_1+p_1'+\cdots+p_M+p_M'=q }}\bigg(\prod_{v=1}^M \frac{\lambda^{p_v}(-\i ta_v)^{p'_v}}{p_v!p'_v!} \bigg) \\ 
        \cdot \sum_{\sigma\in S_q} c_{\sigma, q}R_{\sigma}\big\{[\underbrace{\mc L_M, \ldots,\mc L_M}_{p_M}, \underbrace{\A_{\gamma_M}, \cdots, \mc A_{\gamma_M}}_{p_M'},  \ldots, \underbrace{ \mc L_1,  \ldots, \mc L_1}_{p_1}, \underbrace{\A_{\gamma_1},  \ldots, \mc A_{\gamma_1}}_{p_1'}]\big\}.\label{eq:Z_q}\end{multlined}
    \end{align}
    This matches the BCH formula of $\prod_{v=1}^M\left( e^{\lambda \mc L_{v}} e^{-\i t a_{v} \A_{\gamma_v}} \right)$ (see \citet[Section 3.1]{aftab2024multi}). 
    
    Note that any summand in Eq.~\eq{Z_q} containing $[\mc K_{j_1}, \mc K_{j_2}, \ldots, \mc K_{j_q}]$ for $(j_1, \ldots, j_q)\in [2M]^q$ is a monomial of $\lambda$ and $t$. The degree of $t$ is $|\{k\in [q]\mid j_k \text{ is  odd}|=: q_1$ and the degree of $\lambda$ is $|\{k\in [q]\mid j_k \text{ is  even}\}|=: q_2$. We now bound the coefficient of this monomial. The commutator $[\mc K_{j_1}, \mc K_{j_2}, \ldots, \mc K_{j_q}]$ can only originate from one summand in the summation over $p_v, p'_v$ in Eq.~\eq{Z_q} because the operator counts $(p_1, p_1', \ldots, p_M, p_M')$ are uniquely determined by $(j_1, \ldots, j_q)$. The number of permutations $\sigma\in S_q$ in the second summation in Eq.~\eq{Z_q} that preserve the commutator is \begin{align*}
        \prod_{v=1}^M p_v!p'_v!.
    \end{align*} Thus, the coefficients of the monomial can be upper bounded by \begin{align*}
        \frac{1}{q^2}\prod_{v=1}^M \frac{a_v^{p'_v}}{p_v!p'_v!}\prod_{v=1}^M p_v!p'_v!\big\|[\mc K_{j_1}, \mc K_{j_2}, \ldots, \mc K_{j_q}]\big\| \le  \frac{a_{\max}^{q_1}}{q^2}  \big\|[\mc K_{j_1}, \mc K_{j_2}, \ldots, \mc K_{j_q}]\big\|,
    \end{align*}
    as $|c_{\sigma, q}|\le 1/q^2$. Then, we regroup the summation in Eq.~\eq{Z_q} according to $(j_1, \ldots, j_q)\in [2M]^q$ with the same degree of $t$ and $\lambda$. The resulting polynomial \begin{align*}
        \Omega_{q,t,\lambda}(2M)= \sum_{\substack{q_1, q_2 \ge 0\\q_1+q_2 = q}}\Omega_{q_1, q_2} t^{q_1}\lambda^{q_2}
    \end{align*}
    satisfies \begin{align*}
         \|\Omega_{q_1, q_2}\|\le  \frac{a_{\max}^{q_1}}{q^2}\sum_{(j_1, \ldots, j_q)\in S_{q_1, q_2}} \big\|[\mc K_{j_1}, \mc K_{j_2}, \ldots, \mc K_{j_q}]\big\|.
    \end{align*}
\end{proof}
For many physical Hamiltonians that are $k$-local, the sum of grade-$q$ right-nested commutators scales with $q!$, which causes divergence if we simply bound the remainder by the sum of each term's norm. However, we will show that the remainder of the Magnus expansion can be expressed as the summation of doubly right-nested commutators. The analysis is inspired by \citet{fang2025high}, but as the generators in our case might not be anti-Hermitian, we can not bound the remainder as \cite[Lemma~6]{fang2025high} using the Fourier expansion. Instead, our analysis relies on a detailed analysis of the doubly right-nested commutators in the remainder. In next section, we show that such doubly right-nested commutators admit better scaling than right-nested commutators for $k$-local Hamiltonian so that the remainder of the Magnus expansion can be bounded well. 

Analogous to the $q$-th order Magnus expansion in the preliminaries, we define the $q$-th order approximation of the effective generator as the truncated sum:
\begin{align*}
    \Omega_{(q),t,\lambda}(2M) := \sum_{q'=1}^{q} \Omega_{q',t,\lambda}(2M).
\end{align*}
The corresponding $q$-th order approximate evolution, analogous to $Y_{(q)}(t)$, is then given by
\begin{align*}
    \mc Y_{(q)}(\tau):=\exp\left( \Omega_{(q),t,\lambda}(\tau) \right).
\end{align*}
This approximate evolution is the solution of a differential equation governed by a modified generator $\L_{(q)}(\tau)$, which, analogous to $A_{(q)}(t) $, is given by the series expansion:
\begin{align*}
    \L_{(q)}(\tau) = \sum_{\ell = 0}^{\infty}\frac{1}{(\ell+1)!}\ad_{\Omega_{(q),t,\lambda}(\tau)}^\ell(\dot{\Omega}_{(q),t,\lambda}(\tau)),
\end{align*}
which means \begin{align*}
    \mc Y_{(q)}(\tau) = \exp_{\mc T}\bigg(\int_{0}^{\tau} \mc L_{(q)} (\tau_1) \, \d \tau_1 \bigg).
\end{align*}

We first show that the expansion terms in $\L_{(q)}(\tau)$ can be expressed as doubly right-nested commutators of $\tilde{\mc K}_j$.
\begin{lemma}[Bound on nested commutators of Magnus terms]
\label{lem:doubly-nested-bound-final}
    Define the sum of norms of doubly nested commutators as
    \begin{align*}
    \beta_{\comm,t,\lambda}^{(q_1, \ldots, q_{\ell+1})}(\tau) := \begin{multlined}[t]
        \sum_{j'=1}^{q_{\ell+1}} \sum_{\substack{
            i\in[\ell], j\in[q_i] \\
            i=\ell+1, j\in[q_{\ell+1}-1] \\
            1\le v_{i,j} \le \lceil\tau\rceil
        }} \bigg\| \Big[ 
            [\tilde{\mc K}_{v_{1,1}}, \ldots, \tilde{\mc K}_{v_{1,q_1}}], \ldots, \\ 
             [\tilde{\mc K}_{v_{\ell,1}}, \ldots, \tilde{\mc K}_{v_{\ell,q_\ell}}], 
            [\tilde{\mc K}_{v_{\ell+1,1}}, \ldots, \tilde{\mc K}_{v_{\ell+1,j'-1}},\tilde{\mc K}_{\lceil\tau\rceil},\tilde{\mc K}_{v_{\ell+1,j'}}, \ldots, \tilde{\mc K}_{v_{\ell+1,q_{\ell+1}-1}}]
            \Big] \bigg\|.
    \end{multlined}
\end{align*}
    Then for any sequence of positive integers $q_1, \ldots, q_{\ell+1}$, we have
    \begin{align*}
        \bigg\|\prod_{\ell'=\ell}^1 \ad_{\Omega_{q_{\ell'}, t, \lambda}(\tau)}\cdot\dot\Omega_{q_{\ell+1}, t,\lambda}(\tau)\bigg\| \le \bigg(\prod_{i=1}^{\ell+1} \frac{1}{q_i^2}\bigg) \beta_{\comm}^{(q_1, \ldots, q_{\ell+1})}(\tau).
    \end{align*}
\end{lemma}

\begin{proof}
Let $T_q(t) := \{(\tau_1, \ldots, \tau_q) \mid t \ge \tau_1 \ge \cdots \ge \tau_q \ge 0\}$ be the time-ordered $q$-simplex.
 The term to be bounded is
\begin{align*}
    \mc{X}_{\ell} := [\Omega_{q_1, t, \lambda}(\tau), \ldots, \Omega_{q_\ell, t, \lambda}(\tau), \dot{\Omega}_{q_{\ell+1}, t,\lambda}(\tau)].
\end{align*}
By substituting the integral definitions, using the multilinearity of the commutator, and applying the triangle inequality, we obtain:
\begin{align}
    \|\mc{X}_{\ell}\| 
    &\le \begin{multlined}[t]\sum_{\substack{i = 1 , \ldots, \ell+1\\\sigma_i\in S_{q_{i}}}} \bigg(\prod_{i=1}^{\ell+1} |c_{\sigma_i, q_i}|\bigg) \int_{T_{q_1}(\tau)} \d\tau_{1,1}, \ldots,\d\tau_{1,q_1}\cdots\int_{T_{q_{\ell+1}-1}(\tau)}\d\tau_{\ell+1,1},\ldots,\d\tau_{\ell+1,q_{\ell+1}-1} \\
        \bigg\| \Big[ \mc R_{\sigma_1}\big\{[\mc L(\tau_{1,1}), \ldots, \mc L(\tau_{1,q_1})]\big\}, \ldots, \mc R_{\sigma_\ell}\big\{[\mc L(\tau_{\ell,1}), \ldots, \mc L(\tau_{\ell,q_\ell})]\big\},  \\
        \mc R_{\sigma_{\ell+1}}\big\{[\mc L(\tau), \mc L(\tau_{\ell+1,1}), \ldots, \mc L(\tau_{\ell+1,q_{\ell+1}-1})]\big\} \Big] \bigg\|\nonumber
    \end{multlined}\\
    & \le \begin{multlined}[t]\sum_{\substack{i = 1 , \ldots, \ell+1\\\sigma_i\in S_{q_{i}}}} \bigg(\prod_{i=1}^{\ell+1} \frac{1}{q_i^2} \frac{1}{q_i!}\bigg)q_{\ell+1} \int_{[\tau]^{q_1}} \d\tau_{1,1}, \ldots,\d\tau_{1,q_1}\cdots\int_{[\tau]^{q_{\ell+1}-1}}\d\tau_{\ell+1,1},\ldots,\d\tau_{\ell+1,q_{\ell+1}-1} \\
        \bigg\| \Big[ \mc R_{\sigma_1}\circ \mc T\big\{[\mc L(\tau_{1,1}), \ldots, \mc L(\tau_{1,q_1})]\big\}, \ldots, \mc R_{\sigma_\ell}\circ \mc T\big\{[\mc L(\tau_{\ell,1}), \ldots, \mc L(\tau_{\ell,q_\ell})]\big\},  \\
        \mc R_{\sigma_{\ell+1}}\circ \mc T\big\{[\mc L(\tau), \mc L(\tau_{\ell+1,1}), \ldots, \mc L(\tau_{\ell+1,q_{\ell+1}-1})]\big\} \Big] \bigg\|\nonumber
    \end{multlined}\\
    & \le \begin{multlined}[t]\sum_{\substack{i = 1 , \ldots, \ell+1\\\sigma_i\in S_{q_{i}}}} \bigg(\prod_{i=1}^{\ell+1} \frac{1}{q_i^2} \frac{1}{q_i!}\bigg)q_{\ell+1}
        \sum_{\substack{
            i\in[\ell], j\in[q_i] \\
            i=\ell+1, j\in[q_{\ell+1}-1] \\
            1\le v_{i,j} \le \lceil\tau\rceil
        }} \bigg\| \Big[ \mc R_{\sigma_1}\circ \mc T\big\{[\mc L(v_{1,1}), \ldots, \mc L(v_{1,q_1})]\big\}, \ldots, \\
        \mc R_{\sigma_\ell}\circ \mc T\big\{[\mc L(v_{\ell,1}), \ldots, \mc L(v_{\ell,q_\ell})]\big\},  
        \mc R_{\sigma_{\ell+1}}\circ \mc T\big\{[\mc L(\tau), \mc L(v_{\ell+1,1}), \ldots, \mc L(v_{\ell+1,q_{\ell+1}-1})]\big\} \Big] \bigg\|,\label{eq:bound-Xk}
    \end{multlined} 
\end{align}
where the third line follows from the same integral decomposition argument in the derivation of Eq.~\eq{Z-q-first}. For each $i\in[\ell]$, after applying time-ordering operator $\mc T$, each unordered commutator $[\mc L(v_{i,1}), \ldots, \mc L(v_{i,q_i})]$ is mapped to an ordered commutator \begin{align}
\label{eq:ordered-comm-1}
    [\underbrace{\mc L(\lceil\tau\rceil) \ldots,\mc L(\lceil\tau\rceil)}_{q_{i,\lceil\tau\rceil}},   \ldots, \underbrace{\mc L(1),  \ldots, \mc L(1)}_{q_{i,1}}] = [\underbrace{\tilde{\mc K}_{\lceil\tau\rceil} \ldots,\tilde{\mc K}_{\lceil\tau\rceil}}_{q_{i,\lceil\tau\rceil}},   \ldots, \underbrace{\tilde{\mc K}_{1},  \ldots, \tilde{\mc K}_{1}}_{q_{i,1}}],
\end{align}
where $q_{i,j}$ is the number of $j$ in $v_{i,1}, \ldots, v_{i,q_i}$ and there are \begin{align*}
    \binom{q_i}{q_{i, \lceil\tau\rceil}, \ldots, q_{i,1}} = \frac{q_i!}{q_{i, \lceil\tau\rceil}! \cdots, q_{i,1}!}
\end{align*}
unordered commutators mapped to the same  ordered commutator in Eq.~\eq{ordered-comm-1}. Then, after applying $\mc R_{\sigma_i}$ to each ordered commutator, we regroup the summation according to the index sequence $(u_{i,1}, \ldots, u_{i,q_i})$ of the resulting unordered commutator $[\tilde{\mc K}_{u_{i,1}},\ldots, \tilde{\mc K}_{u_{i,q_i}}]$. Each unordered commutator $[\tilde{\mc K}_{u_{i,1}},\ldots, \tilde{\mc K}_{u_{i,q_i}}]$ is mapped from one ordered commutator $[\underbrace{\tilde{\mc K}_{\lceil\tau\rceil} \ldots,\tilde{\mc K}_{\lceil\tau\rceil}}_{q_{i,\lceil\tau\rceil}},   \ldots, \underbrace{\tilde{\mc K}_{1},  \ldots, \tilde{\mc K}_{1}}_{q_{i,1}}]$ and the number of $\sigma_i\in S_{q_i}$ mapping the ordered commutator to the same unordered commutator is \begin{align*}
    q_{i, \lceil\tau\rceil}! \cdots, q_{i,1}!.
\end{align*}
Therefore, the combinatorial coefficient of $[\tilde{\mc K}_{u_{i,1}},\ldots, \tilde{\mc K}_{u_{i,q_i}}]$ in the final sum is  \begin{align*}
    \frac{q_i!}{q_{i, \lceil\tau\rceil}! \cdots, q_{i,1}!}q_{i, \lceil\tau\rceil}! \cdots, q_{i,1}! = q_i!.
\end{align*}
For $i= \ell+1$, after applying $\mc T$, each $[\mc L(\tau), \mc L(v_{\ell+1,1}), \ldots, \mc L(v_{\ell+1,q_{\ell+1}-1})]$ is mapped to \begin{align}
\label{eq:ordered-comm}
    [\tilde{\mc K}_{\lceil \tau \rceil},\underbrace{\tilde{\mc K}_{\lceil\tau\rceil} \ldots,\tilde{\mc K}_{\lceil\tau\rceil}}_{q_{i,\lceil\tau\rceil}},   \ldots, \underbrace{\tilde{\mc K}_{1},  \ldots, \tilde{\mc K}_{1}}_{q_{i,1}}],
\end{align}
where $q_{i,j}$ is the number of $j$ in $v_{\ell+1, 1},\ldots, v_{\ell+1, q_{\ell+1}-1}$, and there are \begin{align*}
    \binom{q_{\ell+1}-1}{q_{i,\lceil \tau \rceil}, \ldots, q_{i,1}} = \frac{(q_{\ell+1}-1)!}{q_{\ell+1,\lceil \tau \rceil}!, \ldots, q_{\ell+1,1}!}
\end{align*} unordered commutators mapped to the same ordered commutator in Eq.~\eq{ordered-comm}. After applying $\mc R_{\sigma_{\ell+1}}$, we regrouping the summation of the commutators according to the position of the fixed $\tilde{\mc K}_{\lceil \tau \rceil}$ after permutation, $j'$, and the indices of other $\tilde{\mc K}$ operators in a commutator \begin{align*}
    [\tilde{\mc K}_{v_{\ell+1,1}}, \ldots,\tilde{\mc K}_{\lceil\tau\rceil}, \ldots, \tilde{\mc K}_{v_{\ell+1,q_{\ell+1}-1}}].
\end{align*} Each such unordered commutator is mapped from one ordered commutator in Eq.~\eq{ordered-comm} and the number of $\sigma_{\ell+1}\in S_{q_{\ell+1}}$ satisfying $\sigma_{\ell+1}(j')=1$ and $\mc R_{\sigma_{\ell+1}}$ maps the ordered commutator to the same unordered commutator are \begin{align*}
    q_{\ell+1,\lceil \tau \rceil}!, \ldots, q_{\ell+1,1}!.
\end{align*}
Therefore, the combinatorial coefficient of $[\tilde{\mc K}_{v_{\ell+1,1}}, \ldots, \tilde{\mc K}_{v_{\ell+1,j'-1}},\tilde{\mc K}_{\lceil\tau\rceil},\tilde{\mc K}_{v_{\ell+1,j'}}, \ldots, \tilde{\mc K}_{v_{\ell+1,q_{\ell+1}-1}}]$ in the final sum is \begin{align*}
    \frac{(q_{\ell+1}-1)!}{q_{\ell+1,\lceil \tau \rceil}!, \ldots, q_{\ell+1,1}!}q_{\ell+1,\lceil \tau \rceil}!, \ldots, q_{\ell+1,1}! = (q_{\ell+1}-1)!.
\end{align*}
In conclusion, we have \begin{align*}
    \|\mc X_\ell\|&\le \begin{multlined}[t]
        \bigg(\prod_{i=1}^{\ell+1} \frac{1}{q_i^2}\frac{1}{q_i!}q_i!\bigg)q_{\ell+1}/q_{\ell+1}\sum_{j'=1}^{q_{\ell+1}} \sum_{\substack{
            i\in[\ell], j\in[q_i] \\
            i=\ell+1, j\in[q_{\ell+1}-1] \\
            1\le v_{i,j} \le \lceil\tau\rceil
        }} \bigg\| \Big[ 
            [\tilde{\mc K}_{v_{1,1}}, \ldots, \tilde{\mc K}_{v_{1,q_1}}], \ldots, \\ 
             [\tilde{\mc K}_{v_{\ell,1}}, \ldots, \tilde{\mc K}_{v_{\ell,q_\ell}}], 
            [\tilde{\mc K}_{v_{\ell+1,1}}, \ldots, \tilde{\mc K}_{v_{\ell+1,j'-1}},\tilde{\mc K}_{\lceil\tau\rceil},\tilde{\mc K}_{v_{\ell+1,j'}}, \ldots, \tilde{\mc K}_{v_{\ell+1,q_{\ell+1}-1}}]
            \Big] \bigg\|.
    \end{multlined}\\
    &=\begin{multlined}[t]\bigg(\prod_{i=1}^{\ell+1} \frac{1}{q_i^2}\bigg)
        \sum_{j'=1}^{q_{\ell+1}} \sum_{\substack{
            i\in[\ell], j\in[q_i] \\
            i=\ell+1, j\in[q_{\ell+1}-1] \\
            1\le v_{i,j} \le \lceil\tau\rceil
        }} \bigg\| \Big[ 
            [\tilde{\mc K}_{v_{1,1}}, \ldots, \tilde{\mc K}_{v_{1,q_1}}], \ldots, \\ 
             [\tilde{\mc K}_{v_{\ell,1}}, \ldots, \tilde{\mc K}_{v_{\ell,q_\ell}}], 
            [\tilde{\mc K}_{v_{\ell+1,1}}, \ldots, \tilde{\mc K}_{v_{\ell+1,j'-1}},\tilde{\mc K}_{\lceil\tau\rceil},\tilde{\mc K}_{v_{\ell+1,j'}}, \ldots, \tilde{\mc K}_{v_{\ell+1,q_{\ell+1}-1}}]
            \Big] \bigg\|.
    \end{multlined}
\end{align*}
\end{proof}

Define \begin{align}
    &\bar{\beta}_{\comm,t,\lambda}^{(q_1, \ldots, q_{\ell+1})}\nonumber\\
    :=~&\int_{0}^{2M}\beta^{(q_1, \ldots, q_{\ell+1})}_{\comm, t,\lambda}(\tau)\,\d \tau\nonumber\\
    =~&\int_{0}^{2M}\begin{multlined}[t]
        \sum_{j'=1}^{q_{\ell+1}} \sum_{\substack{
            i\in[\ell], j\in[q_i] \\
            i=\ell+1, j\in[q_{\ell+1}-1] \\
            1\le v_{i,j} \le \lceil\tau\rceil
        }} \bigg\| \Big[ 
            [\tilde{\mc K}_{v_{1,1}}, \ldots, \tilde{\mc K}_{v_{1,q_1}}], \ldots, \\ 
             [\tilde{\mc K}_{v_{\ell,1}}, \ldots, \tilde{\mc K}_{v_{\ell,q_\ell}}], 
            [\tilde{\mc K}_{v_{\ell+1,1}}, \ldots, \tilde{\mc K}_{v_{\ell+1,j'-1}},\tilde{\mc K}_{\lceil\tau\rceil},\tilde{\mc K}_{v_{\ell+1,j'}}, \ldots, \tilde{\mc K}_{v_{\ell+1,q_{\ell+1}-1}}]
            \Big] \bigg\|\,\d \tau\nonumber
    \end{multlined} \\
    =~&\begin{multlined}[t]
        \sum_{j'=1}^{q_{\ell+1}}\sum_{v_{\ell+1,j'}=1}^{2M} \sum_{\substack{
            i\in[\ell], j\in[q_i] \\
            i=\ell+1, j\in[q_{\ell+1}]\setminus \{j'\} \\
            1\le v_{i,j} \le v_{\ell+1,j'}
        }} \bigg\| \Big[ 
            [\tilde{\mc K}_{v_{\ell,1}}, \ldots, \tilde{\mc K}_{v_{\ell,q_\ell}}], \ldots, 
            [\tilde{\mc K}_{v_{\ell+1,1}}, \ldots, \tilde{\mc K}_{v_{\ell+1,q_{\ell+1}}}]
            \Big] \bigg\|,\label{eq:def-mueta-bar}
    \end{multlined}
\end{align}
and \begin{align}
    \bar{\beta}_{\comm,q_0,t,\lambda}:=\max\Biggl\{\max_{q>q_0+1}\Biggl(\sum_{\ell=0}^{\infty}\frac{1}{(\ell+1)!}\sum_{\substack{1\le p_1, \ldots, p_{\ell+1}\le q_0\\p_1+\cdots+p_{\ell+1}=q}}\bar{\beta}_{\comm, t,\lambda}^{(p_1, \ldots, p_{\ell+1})}\Biggr)^{1/q}, (\bar{\beta}_{\comm, t,\lambda}^{(q_0+1)})^{1/(q_0+1)}\Biggr\}.
\end{align}
Now, we bound the truncation error of the Magnus expansion in terms of the upper bounds of double right-nested commutators of $\tilde{\mc K}_j$.
\begin{theorem}
    \label{thm:truncation-magnus}
    For any positive integer $q_0$, suppose that \begin{align*}
         \bar{\beta}_{\comm,q_0,t,\lambda}\le 1/2.
    \end{align*}
    Then, the truncation error of the $q_0$-th order Magnus expansion at $\tau = 2M$ is bounded as \begin{align*}
         \|\exp(\Omega_{(q_0), t,\lambda}(2M))-\tilde{\mc P}(t,\lambda)\|\le 2e\bar{\beta}_{\comm,q_0,t,\lambda}^{q_0+1}.
    \end{align*}
\end{theorem}
\begin{proof}
Let $\mc Y(\tau_2, \tau_1) = \exp(-\int_{\tau_1}^{\tau_2}\mc L(\tau')\,\d \tau' )$.
The exponential of $\mc L_{(q_0)}(\tau)$ can be bounded as \begin{align*}
    &\bigg\|\exp_{\mc T}\bigg(\int_{0}^{\tau} \mc L_{(q_0)}(\tau')\,\d \tau'\bigg)\bigg\|  \\
    =~&\bigg\|\mc Y(\tau,0)\sum_{j=0}^{\infty}\int_{0}^{\tau}\d \tau_1\int_0^{\tau_1}\d \tau_2 \cdots \int_{0}^{\tau_{j-1}}\d \tau_j~\prod_{j'=j}^{1}\mc Y(\tau_{j'},0)^{-1}(\mc L_{(q_0)}(\tau_{j'})-\mc L(\tau_{j'}))\mc Y(\tau_{j'},0)\bigg\| \\
    \le ~& \|\mc Y(\tau,0)\|+\bigg\|\sum_{j=1}^{\infty}\int_{0}^{\tau}\d \tau_1\int_0^{\tau_1}\d \tau_2 \cdots \int_{0}^{\tau_{j-1}}\d \tau_j~\prod_{j'=j}^{1}\Big(\mc Y(\tau_{j'-1},\tau_{j'})(\mc L_{(q_0)}(\tau_{j'})-\mc L(\tau_{j'}))\Big)\mc Y(\tau_j,0)\bigg\|\\
    \le ~&1+\sum_{j=1}^{\infty}\int_{0}^{\tau}\d \tau_1\int_0^{\tau_1}\d \tau_2 \cdots \int_{0}^{\tau_{j-1}}\d \tau_j~\prod_{j'=j}^{1}\| \mc L_{(q_0)}(\tau_{j'})-\mc L(\tau_{j'})\|\\
    =~& \exp\bigg(\int_{0}^{\tau}\| \mc L_{(q_0)}(\tau')-\mc L(\tau')\|\,\d \tau'\bigg),
\end{align*}
Then, we can bound the difference between $\exp(\Omega_{(q_0), t,\lambda}(2M))$ and $\tilde{\mc P}(t,\lambda)$ as
\begin{align*}
    &\|\exp(\Omega_{(q_0), t,\lambda}(2M))-\tilde{\mc P}(t,\lambda)\|\\
    =~&\bigg\|\exp_{\mc T}\bigg(\int_{0}^{2M} \mc L_{(q_0)}(\tau)\,\d \tau\bigg)-\exp_{\mc T}\bigg(\int_{0}^{2M} \mc L(\tau)\,\d \tau\bigg)\bigg\| \\
    =~&\bigg\|\int_{0}^{2M} \exp\bigg(\int_{\tau}^{2M} \mc L(s)\,\d s\bigg)(\mc L_{(q_0)}(\tau)-\mc L(\tau)) \exp\bigg(\int_{0}^\tau \mc L_{(q_0)}(s)\,\d s\bigg)\,\d \tau\bigg\| \\
    \le~&\bigg( \int_{0}^{2M}\| \mc L_{(q_0)}(\tau)-\mc L(\tau)\|\,\d \tau\bigg)\exp\bigg(\int_{0}^{2M}\| \mc L_{(q_0)}(\tau)-\mc L(\tau)\|\,\d \tau\bigg).
\end{align*}
where the third line follows from $\|\mc Y(t_2, t_1)\|=1$ for any $t_2 \ge t_1$. Therefore, it is sufficient to bound the difference integral \begin{align*}
    \int_{0}^{2M}\| \mc L_{(q_0)}(\tau)-\mc L(\tau)\|\,\d \tau.
\end{align*}


Now, we upper bound the norm of $\int_{0}^{2M}\|\mc L_{(q_0)}(\tau)-\mc L(\tau)\|\,\d \tau$. By \lem{grade-cond}, we only need to consider terms with a grade greater than $q_0$. By \lem{grade-cond-q+1}, the grade-$(q_0+1)$ term can be bounded by \begin{align*}
    \int_{0}^{2M}\|\dot\Omega_{q_0+1}(\tau)\|\, \d \tau\le\int_{0}^{2M} \beta_{\comm, t,\lambda}^{(q_0+1)}(\tau)\,\d \tau=\bar{\beta}_{\comm, t,\lambda}^{(q_0+1)} \le \bar{\beta}_{\comm,q_0,t,\lambda}^{q_0+1}.
\end{align*}
For any $q>q_0+1$, the grade-$q$ terms can be bounded by \begin{align*}
    &\int_{0}^{2M}\sum_{\ell=0}^{\infty}\frac{1}{(\ell+1)!}\sum_{\substack{1\le p_1, \ldots, p_{\ell+1}\le q_0\\p_1+\cdots+p_{\ell+1}=q}}\beta_{\comm, t,\lambda}^{(p_1, \ldots, p_{\ell+1})}(\tau) \, \d \tau \\
    \le ~&\sum_{\ell=0}^{\infty}\frac{1}{(\ell+1)!}\sum_{\substack{1\le p_1, \ldots, p_{\ell+1}\le q_0\\p_1+\cdots+p_{\ell+1}=q}}\bar{\beta}_{\comm, t,\lambda}^{(p_1, \ldots, p_{\ell+1})} \\
    \le ~&\bar{\beta}_{\comm,q_0,t,\lambda}^q.
\end{align*}
Therefore, the difference integral can be bounded by \begin{align*}
    \int_{0}^{2M}\|\mc L_{(q_0)}(\tau)-\mc L(\tau)\|\,\d \tau&\le \bar{\beta}_{\comm,q_0,t,\lambda}^{q_0+1}+\sum_{q=q_0+2}^{\infty}\bar{\beta}_{\comm,q_0,t,\lambda}^q \\
    &\le \bar{\beta}_{\comm,q_0,t,\lambda}^{q_0+1}+\bar{\beta}_{\comm,q_0,t,\lambda}^{q_0+2}\sum_{j=0}^{\infty}\frac{1}{2^j} \\
    &\le \bar{\beta}_{\comm,q_0,t,\lambda}^{q_0+1}+2\bar{\beta}_{\comm,q_0,t,\lambda}^{q_0+2}\\
    &\le 1/2+2(1/2)^3\le 1
\end{align*}
Therefore, the error between the two exponentials can be bounded as \begin{align*}
\|\exp(\Omega_{(q_0), t,\lambda}(2M))-\tilde{\mc P}(t,\lambda)\|&\le      \bigg( \int_{0}^{2M}\| \mc L_{(q_0)}(\tau)-\mc L(\tau)\|\,\d \tau\bigg)\exp\bigg(\int_{0}^t\| \mc L_{(q_0)}(2M)-\mc L(\tau)\|\,\d \tau\bigg)\\
&\le  e\bar{\beta}_{\comm,q_0,t,\lambda}^{q_0+1}+2e\bar{\beta}_{\comm,q_0,t,\lambda}^{q_0+2} \le 2e\bar{\beta}_{\comm,q_0,t,\lambda}^{q_0+1},
\end{align*}
which concludes the proof.
\end{proof}
Now, we bound the high-order remainder of the difference between $\exp(\Omega_{(q_0), t,\lambda}(2M))$ and ${\mc P}(t,\lambda)$.
Suppose that the Trotter step size is $sT$ for $s\in (0,1)$ and set the physical noise strength to be $\lambda = c(sT)^{p+1}$ for $c > 0$.  Recall that \begin{align*}
    \Omega_{(q_0), sT,c(sT)^{p+1}}(2M)=\sum_{q = 1}^{q_0}\sum_{\substack{q_1,q_2\ge 0\\ q_1+q_2 = q}} \Omega_{q_1, q_2} (sT)^{q_1}(c(sT)^{p+1})^{q_2}.
\end{align*}Define the effective generator of one Trotter step after truncation as \begin{align}
\label{eq:def-effect-generator}
    \G(s)&:= \frac{1}{sT}\Omega_{(q_0), sT,c(sT)^{p+1}}(2M)\\
    &=:\sum_{q=0}^{\infty} \E_{q+1} (sT)^q. 
\end{align}
We bound the superoperators $\E_{q}$ in the following lemma. 
\begin{lemma}
    \label{lem:bound-E}
    Define a constant
    \begin{align}\label{eq:def-mu}
        \mu_{\rm comm} = \max_{q' \ge p+1} \Bigg(\sum_{\substack{1\le q_1+ q_2\le q_0 \\
        q_1+(p+1)q_2 = q'}}\frac{a_{\max}^{q_1}c^{q_2}}{(q_1+q_2)^2} \alpha_{\rm comm}^{(q_1, q_2)}\Bigg)^{1/q'}.
    \end{align}
    Then we have \begin{align}
        \G(s) = -\i \ad_H + \sum_{q=p}^{\infty} \E_{q+1} (sT)^q \label{eq:def-E},
    \end{align} where $\E_{q}$ can be bounded by \begin{align}
        \|\E_q\| \le  \mu_{\rm comm}^q.
    \end{align}
\end{lemma}
\begin{proof}
    If we set $c=0$, by \lem{magnus-bch}, the truncated magnus series $\Omega_{(q_0), sT,0}(2M)$ is the truncated BCH formula of the ideal $p$-th order product formula $\mathcal{P}(sT)$. By the order condition of the product formula (see \citet[Lemma 2]{watson2025exponentially}), we have \begin{align*}
        \frac{1}{sT}\Omega_{(q_0), sT,0}(2M) = -\i \ad_H + O((sT)^p),
    \end{align*}
    which implies \begin{align*}
        \Omega_{q_1, 0} = 0,
    \end{align*}
    for all $1 <q_1 \le p$, and \begin{align*}
        \Omega_{1,0} = - \i \ad_H sT.
    \end{align*} 
    Therefore, for any $c>0$, we have \begin{align*}
        \frac{1}{sT}\Omega_{(q_0), sT,c(sT)^{p+1}}(2M) &= \frac{1}{sT}\Big(\sum_{q_1=1}^{q_0}\Omega_{q_1, 0}(sT)^{q_1} + \sum_{q = 1}^{q_0}\sum_{\substack{q_1\ge 0,q_2\ge 1\\ q_1+q_2 = q}} \Omega_{q_1, q_2} (sT)^{q_1}(c(sT)^{p+1})^{q_2}\Big) \\
        & = -\i \ad_H + \sum_{q_1 = p+1}^{q_0} \Omega_{q_1, 0} (sT)^{q_1-1} + O((sT)^p) \\
        & = -\i \ad_H + \sum_{q = p}^{\infty} \E_{q+1} (sT)^{q}.
    \end{align*}
    For any $q\ge p+1$, we have \begin{align*}
        \|\E_{q}\| &= \bigg\|\sum_{q'=1}^{q_0}\sum_{\substack{q_1, q_2\ge 0,~ q_1+q_2=q' \\
        q_1+(p+1)q_2 = q}} \Omega_{q_1,q_2} c^{q_2} \bigg\|\\
        &\le \sum_{\substack{1\le q_1+ q_2\le q_0 \\
        q_1+(p+1)q_2 = q}}\frac{a_{\max}^{q_1}c^{q_2}}{(q_1+q_2)^2} \alpha_{\rm comm}^{(q_1, q_2)} \\
        &\le  \mu_{\rm comm}^q.
    \end{align*}
    %For $q=p+1$, we have \begin{align*}
    %    \|\E_{p+1}\| &\le c\Big\|\sum_{v=1}^M \L_v\Big\|+\|\Omega_{p+1,0}\| \le c\alpha_{\rm comm}^{(0,1)}+\frac{a_{\max}^{p+1}}{(p+1)^2}\alpha_{\rm comm}^{(p+1,0)}\le \mu_{\rm comm}^{p+1.}
    %\end{align*}
\end{proof}
\begin{lemma}
\label{lem:expansion-exp-G(s)}
    The evolution $e^{\G(s) T}$ admits series expansion in $s$ given by \begin{align*}
        e^{\G(s) T}=e^{-\i \ad_H T}+ \sum_{q=p}^{\infty}\tilde{\E}_{q} s^q,
    \end{align*}
    where $\tilde{\E}_{q}$ are superoperators bounded by \begin{align*}
       \|\tilde{\E}_{q}\|\le   e(2\mu_{\rm comm}T)^{q(1+1/p)},
    \end{align*}
    if $2\mu_{\rm comm}T\ge 1$, and bounded by \begin{align*}
        \|\tilde{\E}_{q}\|\le e(2\mu_{\rm comm}T)^{q},
    \end{align*}
    if $2\mu_{\rm comm}T\le 1$.
\end{lemma}
\begin{proof}
By Dyson series expansion in the interaction picture, we have 
\begin{align*}
    &e^{\G(s) T}-e^{-\i \ad_H T} \\
    = \ & e^{-\i \ad_H T}  \sum_{j=1}^{\infty} \int_{0}^T \d \tau_1 \int_{0}^{\tau_1}\d \tau_2 \cdots \int_{0}^{\tau_{j-1}}\d \tau_j \prod_{j'=j}^{1}\bigg(e^{\i\ad_H \tau_{j'}}\bigg(\sum_{q=p}^{\infty}\E_{q+1}s^qT^q\bigg)e^{-\i \ad_H \tau_{j'}}\bigg) \\
    = \ &  e^{-\i \ad_H T}\sum_{q = p}^{\infty}s^{q} \sum_{j=1}^{\infty} \sum_{\substack{q_1,\ldots,q_j \ge 0 \\ q_1+\cdots+q_j=q}} \int_{0}^T \d \tau_1 \int_{0}^{\tau_1}\d \tau_2 \cdots \int_{0}^{\tau_{j-1}}\d \tau_j \prod_{j'=j}^{1}\bigg(e^{\i\ad_H \tau_{j'}}\E_{q_{j'}+1}e^{-\i \ad_H \tau_{j'}}T^{q_{j'}}\bigg)\\
    =:\,& \sum_{q=p}^{\infty}\tilde{\E}_{q} s^q   
\end{align*}
The superoperator $\tilde{\E}_{q}$ can be bounded by \begin{align*}
    \big\|\tilde{\E}_{q}\big\|&\le \sum_{j=1}^{\infty} \sum_{\substack{q_1,\ldots,q_j \ge 0 \\ q_1+\cdots+q_j=q}} \int_{0}^T \d \tau_1 \int_{0}^{\tau_1}\d \tau_2 \cdots \int_{0}^{\tau_{j-1}}\d \tau_j \prod_{j'=1}^{j}\big(\big\|\E_{q_{j'}+1}\big\|T^{q_{j'}}\big) \\
    &\le  \sum_{j=1}^{\infty} \sum_{\substack{q_1,\ldots,q_j \ge 0 \\ q_1+\cdots+q_j=q}}\frac{T^{q+j}}{j!}\prod_{j'=1}^{j}\mu_{\rm comm}^{q_{j'}+1} \\
    &=\sum_{j=1}^{\lfloor q/p\rfloor }\frac{(\mu_{\rm comm}T)^{q+j}}{j!}\Big(\sum_{\substack{q_1,\ldots,q_j \ge 0 \\ q_1+\cdots+q_j=q}} 1\Big).
\end{align*}
The second summation can be bounded by \begin{align*}
    \sum_{\substack{q_1,\ldots,q_j \ge 0 \\ q_1+\cdots+q_j=q}} 1 &\le \sum_{\substack{ q_1,\ldots,q_j \ge 0 \\ q_1+\cdots+q_j=q}} 1=\binom{q+j-1}{j-1} \le 2^{q+j}.
\end{align*}
Therefore, for $2\mu_{\rm comm}T\ge 1$, we have \begin{align*}
    \big\|\tilde{\E}_{q}\big\|\le \sum_{j=1}^{\lfloor q/p\rfloor }\frac{(2\mu_{\rm comm}T)^{q+j}}{j!} 
    \le (2\mu_{\rm comm}T)^{q+q/p}\sum_{j=1}^{\infty}\frac{1}{j!}\le e(2\mu_{\rm comm}T)^{q(1+1/p)} ,
\end{align*}
and for $2\mu_{\rm comm}T\le 1$, we have \begin{align*}
    \big\|\tilde{\E}_{q}\big\| \le (2\mu_{\rm comm}T)^{q}\sum_{j=0}^{\infty}\frac{ (2\mu_{\rm comm}T)^{j}}{j!}\le (2\mu_{\rm comm}T)^{q} e^{2\mu_{\rm comm}T}\le e(2\mu_{\rm comm}T)^{q}.
\end{align*}
\end{proof}



\section{Applications to $k$-local Hamiltonian}
\label{sec:local}
Suppose that $\ad_H=\sum_{\gamma\in[\Gamma]}\A_\gamma$ and $\mc L_{\tot}:=\sum_{v\in [M]}\mc L_v$ are \emph{$k$-local} on a lattice $\Lambda = [N]$, or specifically \begin{align*}
    |\supp(\A_{\gamma})|, |\supp(\L_v)|\le k, \quad \forall \gamma\in [\Gamma], v\in [M].
\end{align*} For any site $j\in \Lambda$, we further suppose that \begin{align*}
    \sum_{\gamma\in [\Gamma]: j\in \supp(\A_{\gamma})}\|\A_\gamma\|\le g_H, \quad \sum_{v\in [M]: j\in \supp(\L_{v})}\|\L_v\|\le g_L,
\end{align*}
which is called \emph{$g_H$-extensiveness} of $\ad_H$ and $g_L$-extensiveness of $\L_{\rm tot}$.
Recall the unified representation of $\mc A_{\gamma}$ and $\mc L_v$ in Eq.~\eq{def-K}: \begin{align*}
    \mathcal{K}_j = \begin{cases}\mathcal{A}_{\gamma_v} & \text{if } j = 2v-1 \\ \mathcal{L}_v & \text{if } j = 2v \end{cases}, \quad \tilde{\mathcal{K}}_j := \begin{cases}-\i ta_{v}\mathcal{A}_{\gamma_v} & \text{if } j = 2v-1\\ \lambda\mathcal{L}_v & \text{if } j = 2v \end{cases} 
\end{align*}
for $v=1, 2,\ldots, M$. Then, $\sum_{j=1}^{2M} \tilde{\mc K}_j$ is a $k$-local and $(\Upsilon t a_{\max}g_H+\lambda g_L)$-extensive. 

\citet{mizuta2025commutator} provided an upper bound on the sum of right-nested commutators of a $k$-local Hamiltonian, which can be directly extended to the superoperator case. 
\begin{lemma}[{\cite[Lemma~7]{mizuta2025commutator}}]
\label{lem:commu-bound}
    Let $\mc K = \sum_{\gamma=1}^{\Gamma} \mc K_{\gamma}$ be a $k$-local and $g$-extensive Lindbladian on a lattice $\Lambda$ and $\mc A_X$ be a superoperator acting nontrivially on $X\subseteq \Lambda$ with $|X|\le k$, the following inequality on the right-nested commutators holds: \begin{align*}
        \sum_{\gamma_1, \ldots, \gamma_q=1}^{\Gamma} \big\|[\mc K_{\gamma_1}, \ldots, \mc K_{\gamma_{q'}}, \mc A_X, \mc K_{\gamma_{q'+1}}, \ldots, \mc K_{\gamma_q}]\big\|\le q!(2kg)^q\|\mc A_X\|,
    \end{align*}
    for any $q'\in[q]$.
\end{lemma}

\begin{lemma}
\label{lem:commu-bound-sum}
    Let $\mc K_1, \ldots, \mc K_q$ be a sequence of Lindbladian on a size-$N$ lattice $\Lambda$ such that $\mc K_v = \sum_{\gamma=1}^{\Gamma} \mc K_{\gamma}^{(v)}$ is $k$-local and $g_v$-extensive. Then we have 
    \begin{align*}
        \sum_{\gamma_1, \ldots, \gamma_q=1}^{\Gamma} \big\|[\mc K_{\gamma_1}^{(1)}, \ldots, \mc K_{\gamma_q}^{(q)}]\big\|  \le N(q-1)!(2k)^{q-1}\prod_{v=1}^q g_v,
    \end{align*}
\end{lemma}
\begin{proof}
    By summing over the first $q-1$ indices and applying a direct generalization of \lem{commu-bound} to different $\mc K$, we have
    \begin{align*}
        \sum_{\gamma_1, \ldots, \gamma_q=1}^{\Gamma} \big\|[\mc K_{\gamma_1}^{(1)}, \ldots, \mc K_{\gamma_q}^{(q)}]\big\| & \le  (q-1)! \Big(\prod_{v=1}^{q-1} 2kg_v\Big) \sum_{\gamma_q=1}^{\Gamma}\|\mc K_{\gamma_q}^{(q)}\|\\
        &\le (q-1)!\Big(\prod_{v=1}^{q-1} 2k g_v\Big)  Ng_q \\
        &=N(q-1)!(2k)^{q-1}\prod_{v=1}^q g_v.
    \end{align*}
    
\end{proof}
We generalized this result to right-nested commutators of right-nested commutators.
\begin{theorem}[Bound on doubly nested commutators]\label{thm:doubly-nested-bound}
    Let $\mathcal{K} = \sum_{\gamma=1}^{\Gamma} \mathcal{K}_{\gamma}$ be a $k$-local and $g$-extensive Lindbladian on a lattice $\Lambda$, and let $\mathcal{A}_X$ be a superoperator acting nontrivially on $X\subseteq \Lambda$ with $|X|\le k$. For a given sequence of positive integers $q_1, q_2, \ldots, q_\ell$, let $P_{\ell'} = \sum_{j=\ell'}^{\ell} q_j+1$ for $\ell'=1, 2, \ldots, \ell$ and $P_{\ell+1}=1$. Then, the following inequality holds:
    \begin{align*}
        &\sum_{\substack{1 \le i \le \ell,1 \le j \le q_i \\ 1 \le \gamma_{i,j} \le \Gamma}} \bigg\| \Big[ [\mathcal{K}_{\gamma_{1,1}}, \ldots, \mathcal{K}_{\gamma_{1, q_1}}], \ldots, [\mathcal{K}_{\gamma_{\ell-1,1}}, \ldots, \mathcal{K}_{\gamma_{\ell-1, q_{\ell-1}}}], \\
        &\qquad\qquad [\mathcal{K}_{\gamma_{\ell,1}}, \ldots, \mathcal{K}_{\gamma_{\ell, q'}}, \mathcal{A}_X, \mathcal{K}_{\gamma_{\ell, q'+1}}, \ldots, \mathcal{K}_{\gamma_{\ell, q_\ell}}] \Big] \bigg\| \le \bigg(\prod_{i=1}^{\ell}P_{i+1} q_i!(2kg)^{q_i}\bigg)\|\mc A_X\|,
    \end{align*}
    for any $q' \in [q_\ell]$.
\end{theorem}
\begin{proof}
    %Without loss of generality, we assume that each term $\mathcal{K}_{\gamma}$ has a support size of at most $k$. For any $k$-local $\mathcal{K}_{\gamma}$ with larger support, we could decompose it into a sum of terms, each with support size no larger than $k$. By the triangle inequality, the norm of the total commutator would then be bounded by the sum of norms involving these new terms. Since the final bound is independent of the total number of terms $\Gamma$ in the decomposition of $\mc K$, this procedure yields the same result. 
    Define \begin{align*}
        &\mc W_{\gamma_{\ell',1}, \ldots, \gamma_{\ell', q_{\ell'}}, \ldots, \gamma_{\ell, 1}, \ldots, \gamma_{\ell,q_\ell}} =  \Big[ [\mathcal{K}_{\gamma_{\ell',1}}, \ldots, \mathcal{K}_{\gamma_{\ell', q_{\ell'}}}], \ldots, [\mathcal{K}_{\gamma_{\ell-1,1}}, \ldots, \mathcal{K}_{\gamma_{\ell-1, q_{\ell-1}}}], \\
        &\qquad\qquad [\mathcal{K}_{\gamma_{\ell,1}}, \ldots, \mathcal{K}_{\gamma_{\ell, q'}}, \mathcal{A}_X, \mathcal{K}_{\gamma_{\ell, q'+1}}, \ldots, \mathcal{K}_{\gamma_{\ell, q_\ell}}] \Big],
    \end{align*} 
    for $\ell' = 1, 2, \ldots, \ell$ and $q' = 1, 2, \ldots, q_\ell$. We will prove \begin{align} \label{eq:induct-W}
        \sum_{\substack{\ell'\le i\le \ell, 1\le j\le q_i\\ 1\le \gamma_{i,j}\le \Gamma}} \Big\|\mc W_{\gamma_{\ell',1}, \ldots, \gamma_{\ell', q_{\ell'}}, \ldots, \gamma_{\ell, 1}, \ldots, \gamma_{\ell,q_\ell}}\Big\|\le \bigg(\prod_{i=\ell'}^{\ell}P_{i+1} q_i!(2kg)^{q_i}\bigg)\|\mc A_X\|
    \end{align}
    by induction on $\ell'$ from $\ell$ to $1$. For $\ell'=\ell$, it is a direct implication of \lem{commu-bound}. Assume that Eq.~\eq{induct-W} is true for $\ell'$. The LHS of Eq.~\eq{induct-W} for $\ell'-1$ is \begin{align*}
        &\sum_{\substack{\ell'-1\le i\le \ell, 1\le j\le q_i\\ 1\le \gamma_{i,j}\le \Gamma}} \Big\|\mc W_{\gamma_{\ell'-1,1}, \ldots, \gamma_{\ell'-1, q_{\ell'-1}}, \ldots, \gamma_{\ell, 1}, \ldots, \gamma_{\ell,q_\ell}}\Big\|  \\
        =~& \sum_{\substack{\ell'\le i\le \ell, 1\le j\le q_i\\ 1\le \gamma_{i,j}\le \Gamma}} \sum_{\gamma_{\ell'-1,1}\ldots, \gamma_{\ell'-1,q_{\ell'-1}}=1}^{\Gamma}\bigg\|\Big[[\mathcal{K}_{\gamma_{\ell'-1,1}}, \ldots, \mathcal{K}_{\gamma_{\ell'-1, q_{\ell'-1}}}], \mc W_{\gamma_{\ell',1}, \ldots, \gamma_{\ell', q_{\ell'}}, \ldots, \gamma_{\ell, 1}, \ldots, \gamma_{\ell,q_\ell}}\Big]\bigg\|
    \end{align*}
    The summand is nonzero only if one of $\mc K_{\gamma_{\ell'-1,  
    j}}, j\in [q_{\ell'-1}]$ has overlapping support with $\mc W_{\gamma_{\ell',1}, \ldots,\gamma_{\ell,q_\ell}}$. For any fixed $\gamma_{\ell',1}, \ldots, \gamma_{\ell,q_\ell}$, the support size of $\mc W_{\gamma_{\ell',1}, \ldots, \gamma_{\ell,q_\ell}}$ is at most $kP_{\ell'}$. Let $X_{\mc W}$ be the support of $\mc W_{\gamma_{\ell',1}, \ldots,\gamma_{\ell,q_\ell}}$. Then, we have \begin{align*}
        &\sum_{\gamma_{\ell'-1,1}\ldots, \gamma_{\ell'-1,q_{\ell'-1}}=1}^{\Gamma}\bigg\|\Big[[\mathcal{K}_{\gamma_{\ell'-1,1}}, \ldots, \mathcal{K}_{\gamma_{\ell'-1, q_{\ell'-1}}}], \mc W_{\gamma_{\ell',1}, \ldots, \gamma_{\ell', q_{\ell'}}, \ldots, \gamma_{\ell, 1}, \ldots, \gamma_{\ell,q_\ell}}\Big]\bigg\|\\
        \le~&
        \begin{multlined}[t]
            \sum_{j=1}^{q_{\ell'-1}}
            \sum_{\substack{1\le\gamma_{\ell'-1,j}\le \Gamma\\\sup(\mc K_{\gamma_{\ell'-1,j}})\cap X_{\mc W}\neq \emptyset}}
            \sum_{\substack{m=1, \ldots, j-1, j+1, \ldots, q_{\ell'-1}\\1\le\gamma_{\ell'-1,m}\le \Gamma}}
            \big\|[\mathcal{K}_{\gamma_{\ell'-1,1}}, \ldots, \mc K_{\gamma_{\ell'-1,j}}, \ldots, \mathcal{K}_{\gamma_{\ell'-1, q_{\ell'-1}}}]\big\| \\ \cdot2\big\|\mc W_{\gamma_{\ell',1}, \ldots, \gamma_{\ell', q_{\ell'}}, \ldots, \gamma_{\ell, 1}, \ldots, \gamma_{\ell,q_\ell}}\big\|
        \end{multlined} \\
        \le~&\sum_{j=1}^{q_{\ell'-1}}
        \sum_{\substack{1\le\gamma_{\ell'-1,j}\le \Gamma\\\sup(\mc K_{\gamma_{\ell'-1,j}})\cap X_{\mc W}\neq \emptyset}}(q_{\ell'-1}-1)!(2kg)^{q_{\ell'-1}-1}\big\|\mc K_{\gamma_{\ell'-1}, j}\big\|\cdot2\big\|\mc W_{\gamma_{\ell'-1,1}, \ldots, \gamma_{\ell'-1, q_{\ell'-1}}, \ldots, \gamma_{\ell, 1}, \ldots, \gamma_{\ell,q_\ell}}\big\|\\
        \le~& q_{\ell'-1}\cdot 2kP_{\ell'}g\cdot (q_{\ell'-1}-1)!(2kg)^{q_{\ell'-1}-1}\cdot\big\|\mc W_{\gamma_{\ell'-1,1}, \ldots, \gamma_{\ell'-1, q_{\ell'-1}}, \ldots, \gamma_{\ell, 1}, \ldots, \gamma_{\ell,q_\ell}}\big\|\\
        =~& P_{\ell'}(q_{\ell'-1})!(2kg)^{q_{\ell'-1}}\big\|\mc W_{\gamma_{\ell'-1,1}, \ldots, \gamma_{\ell'-1, q_{\ell'-1}}, \ldots, \gamma_{\ell, 1}, \ldots, \gamma_{\ell,q_\ell}}\big\|,
    \end{align*}
    where the third line follows from \lem{commu-bound} and the fourth line follows from $|X_{\mc W}| = kP_{\ell'}$ so we have \begin{align*}
    \sum_{\substack{1\le\gamma_{\ell'-1,j}\le \Gamma\\\mathrm{supp}(\mc K_{\gamma_{\ell'-1,j}})\cap X_{\mc W}\neq \emptyset}}\|\mc K_{\gamma_{\ell'-1}, j}\| \le \sum_{x\in X_{\mc W}}\sum_{\substack{1\le\gamma_{\ell'-1,j}\le \Gamma\\\mathrm{supp}(\mc K_{\gamma_{\ell'-1,j}})\ni x}}\|\mc K_{\gamma_{\ell'-1}, j}\|\le k P_{\ell'} g.
    \end{align*}
    Taking summation over $\gamma_{\ell',1}, \ldots, \gamma_{\ell,q_\ell}$, we have \begin{align*}
        &\sum_{\substack{\ell'-1\le i\le \ell, 1\le j\le q_i\\ 1\le \gamma_{i,j}\le \Gamma}} \Big\|\mc W_{\gamma_{\ell'-1,1}, \ldots, \gamma_{\ell'-1, q_{\ell'-1}}, \ldots, \gamma_{\ell, 1}, \ldots, \gamma_{\ell,q_\ell}}\Big\|  \\
        \le~& P_{\ell'}(q_{\ell'-1})!(2kg)^{q_{\ell'-1}}\sum_{\substack{\ell'\le i\le \ell, 1\le j\le q_i\\ 1\le \gamma_{i,j}\le \Gamma}}
        \Big\|\mc W_{\gamma_{\ell',1}, \ldots, \gamma_{\ell', q_{\ell'}}, \ldots, \gamma_{\ell, 1}, \ldots, \gamma_{\ell,q_\ell}}\Big\|\\
        \le~& \bigg(\prod_{i=\ell'-1}^{\ell}P_{i+1} q_i!(2kg)^{q_i}\bigg)\|\mc A_X\|.
    \end{align*}
\end{proof}
\begin{lemma}
\label{lem:bch-approx} 
    For any positive integer $q_0$ satisfying $4ekq_0(\Upsilon t a_{\max}g_H+\lambda g_L) N^{1/(q_0+1)}\le 1/2$, we have \begin{align*}
        &\Bigg\|\tilde{\P}(t,\lambda)-\exp\bigg( -\i t \ad_H+\lambda\mc L_{\tot}+\sum_{q=2}^{q_0} \Omega_{q,t,\lambda}(2M) \bigg)\Bigg\|\\
        \le~& 2e (4ekq_0(\Upsilon t a_{\max}g_H+\lambda g_L))^{q_0+1}N,
    \end{align*}
    for all $t,\lambda\ge0$.
\end{lemma}
\begin{proof} 
    For any positive integer sequence $q_1, \ldots, q_{\ell+1}$, let \begin{align*}
        P_{\ell'} = \sum_{j=\ell'}^{\ell+1} q_j,
    \end{align*}
    and $q = q_1+\cdots+q_{\ell+1}$. Then, we have \begin{align*}
    &\bar{\beta}_{\comm, t,\lambda}^{(q_1, \ldots, q_{\ell+1})}\\
    =~&\sum_{j'=1}^{q_{\ell+1}}\sum_{v_{\ell+1,j'}=1}^{2M} \sum_{\substack{
            i\in[\ell], j\in[q_i] \\
            i=\ell+1, j\in[q_{\ell+1}]\setminus \{j'\} \\
            1\le v_{i,j} \le v_{\ell+1,j'}
        }} \bigg\| \Big[ 
            [\tilde{\mc K}_{v_{1,1}}, \ldots, \tilde{\mc K}_{v_{1,q_1}}], \ldots, 
            [\tilde{\mc K}_{v_{\ell+1,1}}, \ldots, \tilde{\mc K}_{v_{\ell+1,q_{\ell+1}}}]
            \Big] \bigg\|\\
    \le~& \sum_{j'=1}^{q_{\ell+1}}\sum_{v_{\ell+1,j'}=1}^{2M}\bigg(\prod_{i=1}^{\ell}P_{i+1}q_i!(2k(\Upsilon t a_{\max}g_H+\lambda g_L))^{q_i }\bigg)(q_{\ell+1}-1)!(2k(\Upsilon t a_{\max}g_H+\lambda g_L))^{q_{\ell+1}-1}\|\tilde{\mc K}_{v_{\ell+1, j'}}\| \\
    \le~& \bigg(\prod_{i=1}^{\ell}P_{i+1}q_i!(2k(\Upsilon t a_{\max}g_H+\lambda g_L))^{q_i }\bigg)q_{\ell+1}!(2k(\Upsilon t a_{\max}g_H+\lambda g_L))^{q_{\ell+1}}N\\
    =~& q^{\ell}\bigg(\prod_{i=1}^{\ell+1}q_i!\bigg)(2k)^{q-1}(\Upsilon t a_{\max}g_H+\lambda g_L)^{q}N,
\end{align*}
where the the third line follows from \thm{doubly-nested-bound}, and the fourth line follows from \begin{align*}
    \sum_{v_{\ell+1,j'}=1}^{2M}\|\tilde{\mc K}_{v_{\ell+1, j'}}\|\le \sum_{x\in\Lambda}\sum_{\substack{v = 1, \ldots, 2M \\ \supp(\tilde{\mc{K}}_v)\ni x}} \|\tilde{\mc K}_{v}\|\le(\Upsilon t a_{\max}g_H+\lambda g_L)N.
\end{align*}
For $q_1, \ldots, q_{\ell+1}\le q_0$, we have 
\begin{align*}
    \bar{\beta}_{\comm, t,\lambda}^{(q_1, \ldots, q_{\ell+1})}\le  q^{\ell}q_0^q(2k)^{q-1}(\Upsilon t a_{\max}g_H+\lambda g_L)^{q}N\le q^{\ell+1}(2kq_0(\Upsilon t a_{\max}g_H+\lambda g_L))^{q}N.
\end{align*}
Then, we have 
\begin{align*}
    &\bar{\beta}_{\comm,q_0,t,\lambda}  \\
    =\,&\max\Biggl\{\max_{q>q_0+1}\Biggl(\sum_{\ell=0}^{\infty}\frac{1}{(\ell+1)!}\sum_{\substack{1\le p_1, \ldots, p_{\ell+1}\le q_0\\p_1+\cdots+p_{\ell+1}=q}}\bar{\beta}_{\comm, t,\lambda}^{(p_1, \ldots, p_{\ell+1})}\Biggr)^{1/q}, (\bar{\beta}_{\comm, t,\lambda}^{(q_0+1)})^{1/(q_0+1)}\Biggr\}\\
    \le \,&\max\Biggl\{\max_{q>q_0+1}\Biggl(\sum_{\ell=0}^{\infty}\frac{q^{\ell+1}}{(\ell+1)!}\sum_{\substack{1\le p_1, \ldots, p_{\ell+1}\le q_0\\p_1+\cdots+p_{\ell+1}=q}}(2kq_0(\Upsilon t a_{\max}g_H+\lambda g_L))^{q}N\Biggr)^{1/q}, 2kq_0(\Upsilon t a_{\max}g_H+\lambda g_L)N^{1/(q_0+1)} \Biggr\}\\
    \le \,& \max\Big\{\max_{q>q_0+1}\bigl(e^q(4kq_0(\Upsilon t a_{\max}g_H+\lambda g_L))^{q}N\bigr)^{1/q}, 2kq_0(\Upsilon t a_{\max}g_H+\lambda g_L)N^{1/(q_0+1)}\Big\}\\
    \le \,& \max\Big\{4ekq_0(\Upsilon t a_{\max}g_H+\lambda g_L) N^{1/(q_0+2)}, 2kq_0(\Upsilon t a_{\max}g_H+\lambda g_L)N^{1/(q_0+1)}\Big\} \\
    \le \, &4ekq_0(\Upsilon t a_{\max}g_H+\lambda g_L) N^{1/(q_0+1)}
\end{align*}
where the third line follows from \begin{align*}
    \sum_{\substack{1\le q_1, \ldots, q_{\ell+1}\le q_0\\q_1+\cdots q_{\ell+1}=q}}1\le  \sum_{\substack{1\le q_1, \ldots, q_{\ell+1}\\q_1+\cdots q_{\ell+1}=q}}1=\binom{q-1}{\ell}\le 2^q.
\end{align*}
By \thm{truncation-magnus}, the truncation error is bounded as 
\begin{align*}
\|\exp(\Omega_{(q_0), t,\lambda}(2M))-\tilde{\mc P}(t,\lambda)\|
&\le  2e\bar{\beta}_{\comm,q_0,t,\lambda}^{q_0+1}\\
&\le 2e (4ekq_0(\Upsilon t a_{\max}g_H+\lambda g_L))^{q_0+1}N.
\end{align*}
\end{proof}
Now we substitute $\lambda = c (sT)^{p+1}$ and let $\G(s)$ be the effective generator defined in Eq.~\eq{def-effect-generator}. By \lem{bound-E}, the coefficients of the series expansion of $e^{\G(s)T}$ can be bounded as powers of $\mu_{\rm comm}$ defined in Eq.~\eq{def-mu}. We now provide an upper bound of $\mu_{\rm comm}$. 
\begin{lemma}
\label{lem:bound-mu}
    The parameter $\mu_{\rm comm}$ defined in Eq.~\eq{def-mu} satisfies \begin{align*}
        \mu_{\rm comm} \le q_0 N^{1/(p+1)}(2a_{\max}k \Upsilon g_H+ (2ckg_L)^{1/(p+1)}).
    \end{align*}
\end{lemma}
\begin{proof}
    For the set $S_{q_1, q_2}$ defined in Eq.~\eq{alpha-q-q}, we decompose it according to the position of the odd and even indices.  For any index set $\mathcal{I} \subseteq [q_1+q_2]$, we define the subset $S_{\mathcal{I}}$ as
    \[
    S_{\mathcal{I}} = \Bigl\{ (j_1, \ldots, j_{q_1+q_2}) \in [2M]^{q_1+q_2} \Bigm| 
        j_k \text{ is odd for } k \in \mathcal{I}, \text{ and } 
        j_k \text{ is even for } k \in [q_1+q_2] \setminus \mathcal{I}
    \Bigr\}.
    \]
    We can then express $S_{q_1, q_2}$ as the disjoint union over the collection of all size-$q_1$ subsets of $[q_1+q_2]$:
    \[
    S_{q_1, q_2} = \bigsqcup_{\mathcal{I} \in \binom{[q_1+q_2]}{q_1}} S_{\mathcal{I}}.
    \]
    The commutator bounds are \begin{align*}
    \alpha_{\rm comm}^{(q_1,q_2)} &=  \sum_{(j_1, \ldots, j_{q_1+q_2})\in S_{q_1, q_2}}\big\|[\mc K_{j_1}, \mc K_{j_2},\ldots,\mc K_{j_{q_1+q_2}}]\big\|\\
    &= \sum_{\mc I\in \binom{[q_1+q_2]}{q_1}}\sum_{(j_1, \ldots, j_{q_1+q_2})\in S_{\mc I}}\big\|[\mc K_{j_1}, \mc K_{j_2},\ldots,\mc K_{j_{q_1+q_2}}]\big\|\\
    &\le \sum_{\mc I\in \binom{[q_1+q_2]}{q_1}} N(q_1+q_2-1)!(2k)^{q_1+q_2-1}(\Upsilon g_H)^{q_1}g_L^{q_2}\\
    &= \binom{q_1+q_2}{q_1}N(q_1+q_2-1)!(2k)^{q_1+q_2-1}(\Upsilon g_H)^{q_1}g_L^{q_2}\\
    &\le \binom{q_1+(p+1)q_2}{q_1}N(q_1+q_2)!(2k\Upsilon g_H)^{q_1}((2kg_L)^{1/(p+1)})^{q_2(p+1)},
    \end{align*}
    where the third line follows from \lem{commu-bound-sum}. Then we have  \begin{align*}
        \mu_{\rm comm} &= \max_{q' \ge p+1} \Bigg(\sum_{\substack{1\le q_1+q_2\le q_0 \\
        q_1+(p+1)q_2 = q'}}\frac{a_{\max}^{q_1}c^{q_2}}{(q_1+q_2)^2} \alpha_{\rm comm}^{(q_1, q_2)}\Bigg)^{1/q'} \\
        &\le \max_{q' \ge p+1} \bigg(\sum_{\substack{1\le q_1+q_2\le q_0 \\
        q_1+q_2(p+1) = q'}}\binom{q_1+q_2(p+1)}{q_1}N(q_1+q_2)!(2a_{\max}k\Upsilon g_H)^{q_1}((2ckg_L)^{1/(p+1)})^{q_2(p+1)}\bigg)^{1/q'}\\
        &\le \max_{q' \ge p+1} \bigg(\sum_{q_1=1}^{q_0}\binom{q'}{q_1}Nq_0^{q'}(2a_{\max}k\Upsilon g_H)^{q_1}((2ckg_L)^{1/(p+1)})^{q'-q_1}\bigg)^{1/q'}\\
        &\le q_0N^{1/(p+1)}\max_{q'\ge p+1} \big((2a_{\max}k\Upsilon g_H+(2ckg_L)^{1/(p+1)})^{q'}\big)^{1/q'}\\
        &=q_0N^{1/(p+1)}(2a_{\max}k\Upsilon g_H+(2ckg_L)^{1/(p+1)}).
    \end{align*}
\end{proof}

Now we state our main results for  $k$-local Hamiltonians.  
\begin{theorem}
\label{thm:local-ham}
    Let $p$ be a positive integer, $c$ be a positive number, and $H$ be a $k$-local and $g_H$-extensive Hamiltonian. Suppose that the sum of the noise Lindbladian over the circuit of one $p$-th order product formula is $k$-local and $g_L$-extensive. There is an algorithm that, given any observable $O$ and initial state $\rho_0$,  estimates \begin{align*}
        \tr[O e^{-\i \ad_H T}\rho_0]
    \end{align*}
    to within an error of $\eps\|O\|$ using \begin{align*}
    \widetilde{O}\Big(\frac{( N^{1/(p+1)}\max\{k\Upsilon a_{\max}g_H, (kcg_L)^{1/(p+1)}\}T)^{1+1/p}}{\eps^2}\Big)
    \end{align*}
    Trotter steps in total. 
\end{theorem}
\begin{proof}
    Define the function for extrapolation as \begin{align*}
        f(s)&:=\tr[O(\tilde{\mc P}(sT, c(sT)^{p+1}))^{1/s}\rho_0],
    \end{align*}
    for $s>0$ and $1/s\in \mb Z$.
    Here we consider the scenario where $2\mu_{\rm comm} T \ge 1$. By \lem{expansion-exp-G(s)}, the evolution $e^{\G(s) T}$ admits series expansion in $s$ given by \begin{align*}
        e^{\G(s) T}=e^{-\i \ad_H T}+ \sum_{q=p}^{\infty}\tilde{\E}_{q} s^q,
    \end{align*}
    such that \begin{align*}
        \|\tilde{\mc E}_q\|\le e (2\mu_{\rm comm}T)^{q(1+1/p)}.
    \end{align*}    
    Then the function $f(s)$ can be decomposed as \begin{align*}
        f(s)&= \tr[Oe^{\G(s)T} \rho_0]+\tr[O\bigl((\tilde{\mc P}(sT, c(sT)^{p+1}))^{1/s}-e^{\G(s)T}\bigr)\rho_0] \\
        &= \underbrace{\tr[Oe^{-\i \ad_H T}\rho_0]+\sum_{q=p}^{m-1}\tr[O \tilde{\mc E}_q\rho_0] s^q}_{P_m(s)}+\underbrace{\sum_{q=m}^{\infty}\tr[O\tilde{\mc E}_q\rho_0] + \tr[O(\tilde{\mc P}(sT, c(sT)^{p+1})-e^{\G(s)T})\rho_0]}_{R_m(s)},
    \end{align*}
    where $P_m$ is a degree-$(m-1)$ polynomial in $s$. For any $s$ satisfying \begin{align}
    \label{eq:cond-s}
        4ekq_0\Upsilon sT a_{\max} g_H N^{1/(q_0+1)}\le1/(2e), \quad 4ekq_0c(sT)^{p+1} g_LN^{1/(q_0+1)}\le 1/(2e),\quad s(2\mu_{\rm comm}T)^{1+1/p} \le 1/e,
    \end{align}the remainder $R_m(s)$ is bounded as \begin{align*}
        &|R_m(s)| \\
        \le ~&\sum_{q=m}^{\infty}\|O\|\|\tilde{\mc E}_q\|\|\rho_0\|_1 s^q + \|O\|\Big\|(\tilde{\mc P}(sT, c(sT)^{p+1}))^{1/s}-(e^{\G(s)s})^{1/s}\Big\|\|\rho_0\|_1 \\
        \le ~& \sum_{q=m}^{\infty} e(s(2\mu_{\rm comm} T)^{1+1/p})^{q}\|O\|+\frac{1}{s}\Bigl\|\tilde{\mc P}(sT, c(sT)^{p+1})-e^{\G(s)s}\Bigr\|\|O\|\\
        \le ~&2e^{1-m} \|O\|+ \frac{2e}{s}(4ekq_0(\Upsilon sT a_{\max}g_H+c(sT)^{p+1} g_L))^{q_0+1}N\|O\|\\
        \le ~& 2e^{1-m}\|O\|+2eT\min\{8e^2kq_0\Upsilon  a_{\max} g_H N^{1/(q_0+1)}, (8e^2kq_0c g_LN^{1/(q_0+1)})^{1/(p+1)}\}e^{-(q_0+1)}N\|O\|
    \end{align*}
    where the third line follows from $\tilde{\mc P}(sT, c(sT)^{p+1})$ is a quantum channel and hence $\|\tilde{\mc P}(sT,c(sT)^{p+1})\|=1$, the fourth line follows from \lem{bch-approx} with $t=sT,\lambda=cs^{p+1}$, and the fifth line follows from Eq.~\eq{cond-s} and that the second term is monotonically increasing in $s$. Let $s_1, \ldots, s_m$ and $b_1, \ldots, b_m$ be the extrapolation points and coefficients defined in \lem{richardson-extrapolation} with $t=1$, and then we have $\|\mathbf{b}\|\le C\log m$, $1/s_j\in \mathbb{Z}$, and $1/s_j = r_j = O(\max\{m^3, m^2/s\}/j^2)$. We can set \begin{align*}
        m = O(\log(1/\eps)),
    \end{align*}
    such that \begin{align*}
        2e^{1-m} \le \frac{\eps}{4C\log m}.
    \end{align*}
    Then we can set \begin{align*}
        q_0 = O\Big(\log(NT\min\{k\Upsilon a_{\max}g_H, (kcg_L)^{1/(p+1)}\}/\eps)\Big),
    \end{align*}
    such that \begin{align*}
        2eT\min\{8e^2kq_0\Upsilon  a_{\max} g_H N^{1/(q_0+1)}, (8e^2kq_0c g_LN^{1/(q_0+1)})^{1/(p+1)}\}e^{-(q_0+1)}N\le \frac{\eps}{4C\log m}.
    \end{align*}
    Then by \lem{bound-mu}, the parameter $\mu_{\rm comm}$ satisfies \begin{align*}
        \mu_{\rm comm} & \le q_0 N^{1/(p+1)}(2a_{\max}k \Upsilon g_H+ (2ckg_L)^{1/(p+1)}) \\
        & = O\Big( N^{1/(p+1)}\max\{k\Upsilon a_{\max}g_H, (kcg_L)^{1/(p+1)}\}\log(NT\min\{k\Upsilon a_{\max}g_H, (kcg_L)^{1/(p+1)}\}/\eps)\Big).
    \end{align*}
    Consider the extrapolation function \begin{align*}
        F^{(m)} = \sum_{j=1}^m b_j f(s_j). 
    \end{align*}
    Set \begin{align*}
        f(0) :=\lim_{s\to 0} f(s) = \tr[Oe^{-\i \ad_H T}\rho_0],
    \end{align*} and \begin{align*}
        s &= \min\big\{T^{-1}(8e^2kq_0\Upsilon a_{\max} g_HN^{1/(q_0+1)})^{-1}, T^{-1}(8e^2kq_0 c g_LN^{1/(q_0+1)})^{-1/(p+1)}, e^{-1}(2\mu_{\rm comm}T)^{-1-1/p}\big\} \\
        &= O((\mu_{\rm comm}T)^{-1-1/p}).
    \end{align*}
    Then by \lem{richardson-extrapolation}, we have \begin{align*}
        |F^{(m)}(s)-f(0)|&\le \|\mathbf{b}\|_1 \max_j|R_m(s_j)| \\
        &\le C\log m\Big(\frac{\eps}{4C\log m}+\frac{\eps}{4C\log m}\Big)\|O\|\\
        &\le \frac{\eps }{2}\|O\|.
    \end{align*}
    To estimate $f(s_j)$ with error $\eps\|O\|/(2\|\mathbf{b}\|_1)$ with success probability $1-1/(3m)$, we need to repeatedly run $(\tilde{\mc P}(s_j, cs_j^{p+1}))^{1/s_j}$ $O(\|\mathbf{b}\|_1^2/\eps^2\log m)$ times.  Each application of $(\tilde{\mc P}(s_j, cs_j^{p+1}))^{1/s_j}$ requires $1/s_j$ Trotter steps. Hence the total number of Trotter steps is \begin{align*}
        O\Big(\sum_{j=1}^m r_j\frac{\|\mathbf{b}\|_1^2}{\eps^2}\log m\Big) & = O\Big(\sum_j^{m} \max\{m^3, m^2/s\}\frac{1}{j^2}\frac{\log^3 m}{\eps^2}\Big)\\
        &=O\Big( \max\{\log^3(1/\eps), \log^2(1/\eps)(\mu_{\rm comm}T)^{1+1/p}\}\frac{(\log\log (1/\eps))^3}{\eps^2}\Big) \\
        &= \widetilde{O}\Big(\frac{(\mu_{\rm comm}T)^{1+1/p}}{\eps^2}\Big)=\widetilde{O}\Big(\frac{( N^{1/(p+1)}\max\{k\Upsilon a_{\max}g_H, (kcg_L)^{1/(p+1)}\}T)^{1+1/p}}{\eps^2}\Big)
    \end{align*}
    The smallest noise strength is \begin{align*}
        \min_{j}c(s_jT)^{p+1} &= \Omega\Big(\min_{j}c\Big(\frac{j^2T}{\max\{m^3, m^2/s\}}\Big)^{p+1}\Big) \\
        &= \Omega\Big(c\Big(\frac{T}{\max\{\log^3(1/\eps), \log^2(1/\eps)(\mu_{\rm comm}T)^{1+1/p}\}}\Big)^{p+1}\Big) \\
        &=\Omega\Big(c\min\Big\{\frac{T^{p+1}}{\log^{3p+2}(1/\eps)}, \frac{1}{\log^{2p+2}(1/\eps)\mu_{\rm comm}^{(p+1)^2/p} T^{1+1/p}}\Big\}\Big) \\
        &=\widetilde{\Omega}\Big(c \frac{1}{ N^{1+1/p} \max\{(k\Upsilon a_{\max}g_H)^{(p+1)^2/p}, (kcg_L)^{1+1/p}\} T^{1+1/p}}\Big).
    \end{align*}
    The smallest noise strength is maximized up to a poly-logarithm factor by choosing \begin{align*}
        c = \frac{(k\Upsilon a_{\max}g_H)^{p+1}}{kg_L}, 
    \end{align*}
    which yields \begin{align*}
        \min_{j}c(s_jT)^{p+1} = \widetilde{\Omega}\Big(\frac{1}{ N^{1+1/p} (k\Upsilon a_{\max}g_H)^{1+1/p} kg_L T^{1+1/p}}\Big).
    \end{align*}
    \xz{This reaches the limit of error mitigation under Pauli noise \cite{wang2021noise, cai2021multi}.}
\end{proof}
\section{Applications to geometrically local Hamiltonian}
\label{sec:geo-local}
In this section, we consider a $d$-dimensional lattice $\Lambda \subseteq \mathbb{Z}^d$ and assume that $\ad_H = \sum_{\gamma\in [\Gamma]} \A_{\gamma}$ is geometrical local. 
For any $x\in \Lambda$ and a positive number $R$, define $B_R(x)$ as the $\ell_1$ ball centered on $x$ with radius $R$: \begin{align*}
    B_R(x) = \{ y\in \Lambda:d(x,y) \le R\},
\end{align*}
where \begin{align*}
    d(x,y) = \sum_{i=1}^d  |x_i-y_i|
\end{align*}
is the $\ell_1$ distance between $x = (x_1, \ldots, x_d)$ and $y= (y_1, \ldots, y_d)$. 
Assume that we can group the summand in the decomposition of $\ad_H$ and $\L_{\rm tot}$ as \begin{align*}
    &\ad_H = \sum_{x\in {\Lambda}} \A_x, \quad \A_x = \sum_{\gamma\in \mc I_x} \A_{\gamma} 
\end{align*} where $\{\mc I_x\}_{x\in \Lambda}$ is a partition of $[\Gamma]$, such that \begin{align*}
    &\supp(\A_{\gamma}) \subseteq B_R(x), \quad \forall \gamma\in \mc I_x
\end{align*} for some positive number $R$. We also assume that the observable of interest $O$ is supported on $B_{R_O}(x_O)$ for some positive $R_O$ and $x_O\in \Lambda$. Define \begin{align*}
    J_H = \max_{x\in \Lambda}\sum_{\gamma\in \mc I_x} \|\A_{\gamma}\|.
\end{align*} Define $\Lambda_d = 2^d/d!$ such that \begin{align*}
    |B_R(x)| = \Lambda_dR^d + O(R^{d-1})\le C_d R^d, \quad \forall x\in \Lambda,
\end{align*}
for some $C_d\le 2\Lambda_d$.
Then $\ad_H$ and $\L_{\rm tot}$ are $g_H$ and $g_L$-extensive with \begin{align*}
    g_H &\le \max_{x\in \Lambda} \sum_{y\in B_R(x)}\sum_{\gamma\in \mc I_y} \|\A_{\gamma}\|\le |B_R(x)|J_H  = O(C_d R^d J_H).
\end{align*}
For any $l \ge 0$, define\begin{align*}
    \ad_{H_{(l)}} = \sum_{y\in B_{R_O+l+R}(x_O)} \A_y,
\end{align*}
be the truncated Hamiltonian terms within a ball around the support of $O$. The difference between the time-evolved observables $O$ under $\ad_{H_{(l)}}$ and $ \ad_H$ can be bounded as follows. 
\begin{lemma}[{\cite[Lemma 2]{rao2025stability}}]
    For any $l\ge 0$, we have \begin{align*}
        \left\| e^{\i \ad_H t}O - e^{\i \ad_{H_{(l)}}t}O\right\| \le |\supp(O)|\|O\|\min\{e^{-l/R}(e^{eC_dR^d t}-1), 1\}.
    \end{align*}
\end{lemma}
We can choose \begin{align*}
    l = eC_d R^{d+1}T+ R\log(2|\supp(O)|/\eps),
\end{align*}
such that for any initial state $\rho_0$ \begin{align*}
    \left| \tr[Oe^{-\i \ad_H T}\rho_0]- \tr[Oe^{-\i \ad_{H_{(l)}} T}\rho_0]\right| \le \left\| e^{\i \ad_HT}O - e^{\i \ad_{H_{(l)}}T}O\right\| \le \eps/2.
\end{align*}
Therefore, we can simulate the truncated Hamiltonian $H_{(l)}$ on $B_{R_0+l+2R}(x_O)$,  and the noise terms are restricted to $B_{R_O+l+2R}(x_O)$. Define
\begin{align*}
     \L_{(l), \rm tot} =  \sum_{v\in \in B_{R_O+l+2R}(x_O)} \L_v,
\end{align*}
which includes all possible  noise terms. 
\xz{TBD}
\bibliographystyle{plainnat}
\bibliography{ref}
\end{document}
